\documentclass[manuscript,screen,review]{acmart}

\usepackage{booktabs}
\usepackage{enumitem}
\usepackage{listings}
\usepackage{xspace}

\settopmatter{printacmref=false}
\setcopyright{none}
\renewcommand\footnotetextcopyrightpermission[1]{}

\newcommand{\rasn}{rASN\xspace}
\newcommand{\rdsn}{rDSN\xspace}
\newcommand{\codepilot}{CodePilot\xspace}
\newcommand{\cpp}{C\nolinebreak\hspace{-.05em}\raisebox{.4ex}{\tiny\bf ++}\xspace}

\lstdefinestyle{rasnstyle}{
  basicstyle=\ttfamily\small,
  breaklines=true,
  columns=fullflexible,
  frame=single,
  keepspaces=true,
  showstringspaces=false
}
\lstset{style=rasnstyle}

\title{Robust Agent System Nucleus: A Distributed-Systems-Inspired Runtime for Robust AI Agent Systems}

\author{rASN Project}
\affiliation{%
  \institution{Robust Agent System Nucleus Prototype}
  \city{Beijing}
  \country{China}}
\email{technical-report@rasn.local}

\begin{abstract}
Modern AI agents can solve complex software-engineering tasks, but their
execution behavior is often difficult to reproduce, diagnose, and control.
Failures are frequently attributed to model quality, while the surrounding
runtime exposes weak abstractions for nondeterminism, hidden dependencies,
state, workflow execution, and observability. This report presents the design
and current implementation of Robust Agent System Nucleus (\rasn), a prototype
\cpp framework for building robust agent systems using principles from robust
distributed systems and systems tooling. \rasn is built on top of Robust
Distributed System Nucleus (\rdsn), reusing its service-app model, typed RPC
task codes, serverlet/clientlet abstractions, configuration, timers, locks, and
logging. \rasn models agents as explicit service components, represents agent
work as versioned task messages, records nondeterministic choices as runtime
events, persists state through a named state service, and compiles declarative
workflows into executable agent graphs. The first application is \codepilot, a
coding CLI inspired by contemporary coding agents. \codepilot routes model and
tool operations through the generic \rasn coordinator, supports simulator and
HTTP-compatible LLM providers, gates local side effects through policy, and can
run either inline or as a typed \rdsn service graph. We describe the
architecture, implementation, expected evaluation methodology, threats to
validity, and open research/product directions. The current report is intended
both as a technical artifact for the prototype and as a living research report
that should be updated as \rasn and \codepilot evolve.
\end{abstract}

\begin{CCSXML}
<ccs2012>
 <concept>
  <concept_id>10011007.10011006.10011008</concept_id>
  <concept_desc>Software and its engineering~Software architectures</concept_desc>
  <concept_significance>500</concept_significance>
 </concept>
 <concept>
  <concept_id>10010520.10010521.10010537</concept_id>
  <concept_desc>Computer systems organization~Distributed architectures</concept_desc>
  <concept_significance>500</concept_significance>
 </concept>
 <concept>
  <concept_id>10002951.10003260.10003261</concept_id>
  <concept_desc>Information systems~Data management systems</concept_desc>
  <concept_significance>300</concept_significance>
 </concept>
</ccs2012>
\end{CCSXML}

\ccsdesc[500]{Software and its engineering~Software architectures}
\ccsdesc[500]{Computer systems organization~Distributed architectures}
\ccsdesc[300]{Information systems~Data management systems}

\keywords{AI agents, distributed systems, runtime systems, observability, replay, workflow orchestration, coding agents}

\begin{document}

\maketitle

\section{Introduction}

AI agent systems increasingly combine large language models (LLMs), tool
execution, persistent context, and multi-step orchestration to automate
software-engineering tasks. Coding assistants such as Codex-style command-line
interfaces, GitHub Copilot CLI, Claude Code, OpenCode, and similar systems
demonstrate that LLM-driven agents can perform useful task-level work. However,
their runtime behavior often remains fragile. A single user request can involve
multiple model calls, tool calls, filesystem effects, hidden provider
configuration, retries, timeouts, and state updates. When the result is wrong,
slow, expensive, or unsafe, operators often lack a structured execution model
that explains what happened and how to reproduce it.

This report argues that many robustness problems in agent systems should be
treated as systems problems rather than only as model problems. From this
perspective, an agent application resembles a distributed service system: it
contains multiple components with explicit identities, independent failure
modes, asynchronous boundaries, state ownership, service discovery, request
deadlines, retry policies, side effects, and observability requirements.
Distributed systems research and infrastructure have long studied how to expose
these concerns through runtime abstractions and operational tooling. Robust
Distributed System Nucleus (\rdsn)~\cite{rdsn} is one such framework: it provides a
microkernel-style foundation for building robust distributed services with
pluggable service applications, typed RPC, local runtimes, configuration,
logging, task execution, and stateful components.

Robust Agent System Nucleus (\rasn) explores the hypothesis that an agent
runtime can become more reliable and diagnosable by reusing these distributed
systems principles. The current \rasn prototype is a \cpp plugin under the
\rdsn source tree. It defines an agent service graph consisting of a registry,
coordinator, model agent, tool agent, state service, workflow service,
observability service, and application gateway. Each component can execute
inline for local prototyping or as a \rdsn service app exposing typed RPC
handlers. The runtime records nondeterministic model and simulator decisions,
captures failure events, stores workflow state and leases through a namespaced
state service, and supports replay-oriented diagnosis.

The first application on top of \rasn is \codepilot, a coding CLI. \codepilot
implements provider selection, local tools, reusable coding skills, workflow
commands, state inspection, registry inspection, observability commands, and
self-tests. Importantly, \codepilot does not call a model or local tool
directly. It converts user intents into generic \texttt{agent\_request}
messages and routes them through the \rasn coordinator. This separation keeps
coding-agent features outside the generic runtime and allows future
applications to reuse the same service graph. Tool providers expose structured
descriptors with side-effect classes and argument schemas, so orchestration and
UI layers can inspect tool surfaces from typed metadata instead of scraping a
prose-only command list.

This report makes the following contributions:

\begin{itemize}
  \item We formulate a design philosophy for robust agent runtimes based on
  distributed service principles: explicit component boundaries, typed
  messages, durable state, observable nondeterminism, bounded retries,
  cancellation, and workflow leases.
  \item We describe the current \rasn architecture and implementation on top of
  \rdsn, including service-mode RPCs, state recovery, workflow execution,
  observability with \rdsn \texttt{perf\_counter} metrics, provider adapters,
  policy-bounded tools, and lifecycle hardening.
  \item We present \codepilot as an application case study and explain how a
  coding CLI can be implemented as an adapter over generic agent services.
  \item We define an experimental design for future evaluation, including
  environment setup, metrics, workloads, and comparison targets. This section is
  intentionally forward-looking because \rasn is still a prototype.
  \item We discuss threats to validity, practicality, generality,
  extensibility, and the remaining work required to turn \rasn into a solid
  product-quality platform.
\end{itemize}

\section{Background and Motivation}

\subsection{Robust Distributed System Nucleus}

\rdsn was designed to accelerate the construction of robust distributed
systems. Its microkernel architecture separates applications, distributed
frameworks, DevOps tooling, and local runtime providers. Developers structure
logic as service applications, define RPC task codes, register server-side
handlers, create client wrappers, and use shared runtime services such as
configuration, logging, task scheduling, and filesystem utilities. The \rasn
prototype follows patterns from existing \rdsn applications such as echo-style
RPC services and stateful key-value services.

The important design lesson for \rasn is that robust behavior is not an
afterthought. It emerges from explicit runtime boundaries. Service identities
are named; RPC operations are typed; state is owned by a service; timers and
background tasks are explicit; configuration is visible; failures are returned
as structured errors; and operational tests can start the entire service graph
with a configuration file. These are precisely the properties that many agent
systems lack.

\subsection{AI Agent Runtime Problems}

An AI agent system combines deterministic code with nondeterministic model
calls and side-effecting tools. Four classes of problems recur in practice:

\begin{enumerate}
  \item \textbf{Uncontrolled nondeterminism.} Model sampling, provider
  responses, tool output, retry timing, and dynamic context can change between
  executions. Without a runtime log, reproduction is weak.
  \item \textbf{Hidden dependencies.} Agents often depend on environment
  variables, provider credentials, local files, shell commands, vector stores,
  and external services. Failure boundaries are not first-class.
  \item \textbf{Weak state management.} Conversation state, workflow progress,
  tool artifacts, and checkpoints may be kept in memory or ad hoc files.
  Recovery semantics are unclear.
  \item \textbf{Limited observability and control.} Operators need to inspect
  routes, failures, retries, deadlines, cancellation, and state. Many current
  tools expose only a transcript.
\end{enumerate}

\rasn treats these as runtime design requirements. The goal is not to replace
LLM research, but to provide a systems substrate in which model behavior,
tooling, state, and orchestration can be made explicit and testable.

\subsection{Coding Agents as a First Application}

Coding CLIs are an appropriate first application because they stress several
runtime concerns simultaneously. They need LLM provider adapters, local project
inspection, filesystem read/search tools, optional write and shell tools,
workflow decomposition, user-facing status, and robust diagnostics when a build
or tool call fails. They also require conservative policy defaults: reads are
safe, writes and shell commands must be opt-in, and large outputs need
bounded artifact handling.

\section{Design Philosophy}

\rasn is guided by six design principles.
The principles are not merely presentation choices; they are constraints that
shape where state is stored, where errors are surfaced, and which parts of the
system are allowed to perform side effects. The main design decision is to move
agent robustness concerns out of prompts and into runtime-visible mechanisms.
For example, cancellation is represented as request/runtime state, not as an
instruction to the model; workflow progress is represented as state records, not
as a transcript paragraph; and model/tool calls are represented as typed tasks,
not as free-form strings interpreted differently by each application.

This philosophy follows the same engineering pattern as \rdsn. \rdsn does not
make a distributed system robust by providing a single monolithic library; it
provides service lifecycle, RPC, task codes, configuration, failure handling,
and tooling so applications can build robust protocols. \rasn applies the same
idea to agent systems. It provides a nucleus of runtime abstractions---agents,
capabilities, workflows, state, policy, and observability---while leaving
domain-specific prompts and interfaces to applications such as \codepilot.

\subsection{Principle-to-Mechanism Mapping}

Table~\ref{tab:principle-mechanism} maps each principle to concrete mechanisms
in the current prototype. This mapping is useful because it makes the design
falsifiable: if a mechanism is absent, the corresponding principle is not yet
implemented.

\begin{table*}[t]
\caption{Design principles and concrete mechanisms.}
\label{tab:principle-mechanism}
\centering
\begin{tabular}{p{0.18\textwidth}p{0.36\textwidth}p{0.34\textwidth}}
\toprule
Principle & Mechanisms & Current limitations \\
\midrule
Service components & \rdsn service apps, serverlets/clientlets, descriptors,
capabilities, health, readiness probes, URI/host endpoint configuration,
generated client surfaces & Inline mode keeps a same-process fallback for local
prototyping; multi-process deployment testing remains future work \\
Typed messages & Versioned structs, binary serialization, schema validation,
structured errors, runtime schema manifest, JSON/IDL export, generated
C++/TypeScript/Python contract stubs, generated C++/TypeScript/Python RPC
clients & Product packaging and transport implementations remain future SDK
work \\
Explicit nondeterminism & Trace events, simulator choices, replay-load hooks,
model/tool result replay, scheduler-order checks, filesystem snapshots,
external-effect ledger entries, environment-provider time and randomness, trace
ids & Full external-environment virtualization remains future work \\
State boundary & \texttt{rasn.state}, CAS, checkpoints, journals, leases,
artifact/snapshot indexes, optional local replica mirror and NFS import &
Backing store is not quorum-replicated \\
Policy boundaries & Tool classification, workspace-root sandbox, opt-in
writes/shell, human-approval labels, shell command allowlists, bounded output,
artifact spill, redaction boundary, model credential handles, shell timeout/job
containment, configurable container command wrapper & Hardened container
orchestration remains future deployment work \\
Iterative hardening & Unit tests, self-tests, service smoke, documented phases,
report updates, \codepilot evaluation harness & User-study evaluation remains
future empirical work \\
Operational observability & Structured JSONL traces, failure classification,
diagnosis helpers, and \rdsn \texttt{perf\_counter} metrics (counters and
latency percentiles) exported as text/Prometheus/JSON via \texttt{observe
metrics} and a \texttt{command\_manager} command & No bundled scrape gateway,
retention store, or dashboards \\
Resilience & Bounded coordinator retry, request deadlines, per-provider
model circuit/admission/rate guards, per-tool admission/rate gates, and
per-remote-agent RPC circuit/admission/rate guards, all surfaced through
\texttt{perf\_counter} metrics and a \texttt{rasn.resilience} command & The
limits are per-dependency rather than a single process-wide budget, and govern
request count rather than token/cost budgets \\
\bottomrule
\end{tabular}
\end{table*}

\subsection{Agents as Service Components}

Each agent is modeled as a service component with identity, lifecycle,
capabilities, health, configuration, and endpoint metadata. The built-in graph
contains:

\begin{table}[t]
\caption{\rasn service roles in the current prototype.}
\label{tab:roles}
\centering
\begin{tabular}{ll}
\toprule
Role & Responsibility \\
\midrule
\texttt{rasn.registry} & Agent descriptors and capability lookup \\
\texttt{rasn.coordinator} & Routing, retries, and orchestration boundary \\
\texttt{rasn.llm.agent} & LLM provider isolation and model gateway \\
\texttt{rasn.tool.agent} & Local tool isolation and policy enforcement \\
\texttt{rasn.state} & Namespaced state, checkpoint, journal, CAS \\
\texttt{rasn.workflow} & Workflow validation, compilation, execution \\
\texttt{rasn.observability} & Event, failure, replay, and snapshot queries \\
\texttt{rasn.codepilot} & Coding CLI gateway application \\
\bottomrule
\end{tabular}
\end{table}

The same services can execute inline or through \rdsn RPC. Inline mode supports
fast local development; service mode exercises typed RPC task codes and service
lifecycle behavior.

The service-component view also defines a boundary of accountability. A model
agent is accountable for provider health, provider configuration, request
timeouts, and response normalization; it is not accountable for deciding whether
a shell command is safe. A tool agent is accountable for executing local tools
under policy; it is not accountable for choosing which workflow node should run
next. A workflow service is accountable for dependency ordering, state
transitions, and lease ownership; it is not accountable for parsing CodePilot's
interactive commands. These boundaries make it possible to replace one module
without changing the others. For instance, a future replicated state service can
replace the local state store if it preserves put/get/query/checkpoint/CAS
semantics.

This differs from a common coding-agent architecture in which a CLI owns the
prompt loop, tool list, conversation memory, local filesystem mutations, and
provider calls in one process. Such an architecture is easy to prototype but
hard to debug after failures because the key decisions are implicit in the
transcript. \rasn deliberately pays more upfront engineering cost to obtain
separate lifecycle, routing, state, policy, and observability surfaces.

\subsection{Typed Messages Rather Than String Envelopes}

The public runtime protocol centers on versioned request, response, descriptor,
and error structures. These messages carry request IDs, trace IDs, workflow IDs,
capability names, inputs, context entries, timeout budgets, retry budgets,
policy labels, and replay mode. All top-level messages are serialized through
\rdsn binary readers and writers. The design avoids burying routing or state
semantics inside prompt text.

Typed messages provide three robustness benefits. First, they make validation
local: a missing capability or unsupported schema version can be rejected before
any provider observes the request. Second, they make cross-component evolution
explicit: a new field must either have a backward-compatible default or be
guarded by a schema-version check. Third, they give observability a stable
schema. A trace viewer can search for \texttt{request\_id}, \texttt{trace\_id},
or \texttt{capability} without understanding the natural-language prompt.

The prototype also exposes a runtime schema manifest. The manifest lists the
current fields, types, requiredness, version, and purpose for core contracts
such as \texttt{agent\_request}, \texttt{agent\_response},
\texttt{policy\_request}, \texttt{model\_provider\_descriptor},
\texttt{runtime\_event}, \texttt{state\_record}, and \texttt{workflow\_node}.
CodePilot exposes this catalog through
\texttt{codepilot schema}. The same manifest can now be exported as JSON
(\texttt{codepilot schema json}), as a compact rASN IDL
(\texttt{codepilot schema idl}), as a self-contained generated C++ contract
header (\texttt{codepilot schema cpp}), as TypeScript interfaces
(\texttt{codepilot schema ts}), or as Python dataclasses
(\texttt{codepilot schema py}). A separate RPC operation manifest enumerates
service names, task codes, request types, and response types, and
\texttt{codepilot schema clients-cpp}, \texttt{clients-ts}, and
\texttt{clients-py} emit executable client wrappers for agent, registry, state,
workflow, model, and observability RPCs. The C++ client output uses \rdsn
\texttt{clientlet}; the TypeScript and Python outputs target a small transport
interface so a host application can provide HTTP, local IPC, or \rdsn-backed
bridges. These generators cover core agent, registry, tool, policy,
observability, state, and workflow records, including nested types referenced by
top-level requests. Unit tests check that the manifest contains these contracts,
that generated C++, TypeScript, and Python surfaces include typed vectors/lists
and schema-version fields, and that generated client surfaces include
representative state, workflow, and observability task codes. This narrows the
SDK gap to product packaging and concrete multi-language transports rather than
schema or client surface generation.

\rasn intentionally does not use a generic JSON blob as the core runtime
envelope. JSON is convenient for provider APIs and configuration, but it is too
weak as an internal systems boundary because required fields, versioning,
timeouts, side-effect labels, and error categories become conventions rather
than compiler-visible interfaces. The prototype still uses strings for payloads
where the value is inherently provider-specific, but routing and control data
are kept in typed fields.

\subsection{Nondeterminism as Runtime Data}

\rasn records runtime events such as \texttt{task.begin},
\texttt{llm.request}, \texttt{llm.response}, \texttt{nondeterminism},
\texttt{failure}, \texttt{retry}, and \texttt{task.finish}. The simulator
provider records its random choice as a named nondeterministic value. Replay can
load a prior trace and reuse recorded nondeterministic values. This is not yet a
complete deterministic replay system, but it establishes a first-class place
for nondeterminism in the runtime model.

The design distinguishes several sources of nondeterminism: model sampling,
provider availability, tool output that depends on the local environment,
workflow scheduling among simultaneously ready nodes, retry timing, and external
state recovery. \rasn does not try to eliminate these sources. Instead, it
captures enough information to make them inspectable. For example, a model
request records provider, model, prompt/context metadata, timeout budget, and
result; the simulator additionally records the selected pseudo-random output;
workflow execution records the node order and terminal state; and state records
capture durable decisions such as leases and cancellations.

In-process timing and randomness are themselves drawn from rDSN's pluggable
environment provider (\texttt{dsn\_now\_ms} and \texttt{dsn\_random64}, the same
primitive rDSN uses to mint RPC trace ids) rather than from private
wall-clock-seeded generators. Because that provider is an overridable rDSN
component whose thread-local generator can be seeded by tooling, emulator and
replay tools can virtualize or reproduce these inputs without changing \rasn
code. This makes time and randomness first-class, controllable environment
inputs instead of hidden sources of divergence.

This choice reflects a practical research hypothesis: robust agent systems do
not require every execution to be deterministic, but they do require
nondeterministic decisions to be named, logged, and, when possible, replayable.
The current prototype therefore treats replay as an incremental capability. The
first milestone is trace-guided diagnosis; later milestones can add
deterministic provider simulation, cached model responses, and scheduler replay.

\subsection{State Behind an Explicit Service Boundary}

Workflow runs, workflow node status, leases, observability snapshots, and
oversized tool-output metadata are written through the state service boundary.
The state service supports namespaced keys, checkpoint files, append-only
journals, recovery, and conditional writes. This design prevents hidden
in-process state writes from bypassing service-mode behavior.

State is divided into two categories. \emph{Coordination state} affects future
runtime decisions: workflow run records, node records, and leases determine
whether work can resume, cancel, or complete. \emph{Diagnostic state} improves
operator visibility: artifact indexes and observability snapshots point to
larger files or traces. Both categories cross \texttt{rasn.state}, but they have
different correctness requirements. Coordination state uses sequence checks when
overwrites could violate ownership or terminal-state invariants; diagnostic
state can often be append-like or best-effort as long as errors are visible.

The state boundary is also a preparation point for production. The current store
is local and simple, but all higher-level modules depend on the state-service
interface rather than a concrete map. This means future work can substitute a
replicated key-value store, a database-backed store, or an \rdsn SKV-backed
store while preserving the workflow and policy modules' contracts.

\subsection{Safety Through Policy and Boundaries}

The tool agent evaluates every CodePilot tool invocation through a policy
manager. Read-only operations are allowed only for existing paths; write and
shell operations are disabled by default; network-class operations are also
explicitly gated. Tool output is bounded, and oversized output is spilled to an
artifact file with metadata stored in state. These mechanisms are intentionally
simple but demonstrate the runtime shape needed for stronger policy engines.

The policy philosophy is conservative: observation should be easy, mutation
should be explicit, and unbounded effects should be converted into bounded
runtime artifacts. This is why read/list/search are treated differently from
write/replace/shell. It is also why a large tool response is not simply pushed
into the next model prompt. The runtime creates a bounded preview and a durable
artifact reference, giving the user and future tools a place to inspect the full
output without losing control of prompt size.

The current policy manager is not a process sandbox. It is a policy boundary:
all local side effects pass through a single module where repository scoping,
approval gates, command allowlists, secret redaction, and audit logging can be
centralized. The design therefore favors a small set of explicit tools over a
single unrestricted shell primitive.

\subsection{Product Readiness as Iterative Hardening}

The prototype has evolved through many hardening rounds. Recent examples
include reference-counted service-graph lifecycle, lazy \codepilot graph
ownership, full service readiness probes, workflow execution leases, lease
renewal, state-bound artifact writes, cooperative agent cancellation,
cancellation tombstone retention, state-recovery error propagation, and
restart-safe workflow cancellation. This
iterative pattern is important: each refinement turns a previously implicit
assumption into an explicit runtime invariant and test.

The hardening process follows a repeating pattern:

\begin{enumerate}
  \item identify an implicit assumption, such as ``state lookup failure means the
  workflow is missing'';
  \item turn the assumption into a named invariant, such as error transparency;
  \item move the invariant into a runtime boundary, such as state recovery or
  workflow cancellation;
  \item add a regression test or self-test path; and
  \item update the design documentation and this report.
\end{enumerate}

This process is part of the research contribution. It treats robustness as an
engineering discipline for agent runtimes, not as a property that emerges only
from using a stronger model. The prototype is still incomplete, but its history
shows how distributed-systems-style hardening can be applied incrementally to an
agent platform.

\section{Runtime Model and Robustness Invariants}

\rasn treats an agent system as a set of typed, stateful service components that
communicate through explicit requests and persist important runtime decisions.
This section defines the current prototype's runtime entities and the invariants
that guide the implementation. The model is intentionally practical rather than
fully formal: each entity below corresponds to concrete \cpp structs, RPC task
codes, state records, or log events in the implementation.

\subsection{Failure Model}

The current runtime model assumes crash-stop and omission-style failures at
component boundaries rather than Byzantine behavior. A provider may time out,
return malformed data, or report an API error. A tool may fail, produce bounded
or oversized output, or be denied by policy. A service-mode RPC may fail because
the endpoint is not ready, the task times out, or serialization rejects the
message. The state service may report a missing key, a conditional-write
conflict, or a persistence/recovery error. The registry may contain stale dynamic
descriptors until heartbeat leases expire and sweeps run.

These failures are represented explicitly. \rasn avoids three patterns that make
agent systems difficult to debug: treating infrastructure failures as empty
model output, treating state errors as missing state, and treating side-effect
timeouts as safe-to-retry failures. The coordinator, workflow service, and
CodePilot readiness checks all propagate boundary-specific errors so the caller
can distinguish a provider failure from a registry failure, state conflict, or
policy denial.

\subsection{Consistency Model}

\rasn does not currently provide distributed consensus or replicated state. Its
consistency model is therefore scoped to a single state-service instance, but it
still matters because many robustness bugs are local race bugs. Within the state
store, each successful write receives a monotonic sequence number. Conditional
puts compare the caller's expected sequence with the current record sequence.
Workflow leases and terminal workflow transitions rely on this compare-and-set
discipline.

The model can be summarized as follows:

\begin{itemize}
  \item reads observe the latest record in the local state-service instance;
  \item unconditional writes replace the record and advance the sequence;
  \item create-only writes succeed only when the key is absent;
  \item check-sequence writes succeed only when the current sequence matches the
  expected sequence;
  \item checkpoints, journals, local replica mirrors, and NFS imports are recovery
  artifacts, not consensus protocols; and
  \item service-mode callers should treat state errors as hard errors unless the
  error is the explicit missing-key response for the requested key.
\end{itemize}

This is intentionally weaker than a production replicated store, but it is
strong enough to model ownership, terminal-state protection, restart-safe query,
restart-safe cancellation, and operator recovery from a secondary local copy. The
report's evaluation plan therefore separates ``runtime invariant correctness''
from future ``replicated storage availability.''

\subsection{Identity and Correlation}

Robustness depends on stable identifiers. \rasn uses distinct identifiers for
different purposes: \texttt{request\_id} for duplicate detection and
cancellation, \texttt{trace\_id} for event correlation, \texttt{task\_id} for
logical work, \texttt{run\_id} for workflow executions, \texttt{node\_id} for
workflow tasks, \texttt{agent\_id} for service instances, and state keys for
durable records. Keeping these identifiers separate avoids a common failure mode
in agent systems where conversation ids, tool-call ids, and workflow ids are
interchangeably embedded in text.

The identifiers form a hierarchy during workflow execution: a workflow run owns
node ids; each node dispatches model/tool requests with request ids; all these
requests share or derive trace context; state records preserve run and node ids;
observability events preserve trace ids and task ids. This makes it possible to
answer questions such as ``which provider call produced this node output?'' or
``which state record made this cancellation terminal?'' without reconstructing
causality from natural language.

\subsection{Runtime Entities}

\begin{description}
  \item[Agent descriptor.] A descriptor names an agent instance, role, app name,
  endpoint, health state, version, supported capabilities, and optional lease
  metadata. Descriptors are stored in the registry and are the unit of routing.
  \item[Capability.] A capability is a typed operation exposed by an agent, such
  as \texttt{model.complete}, \texttt{tool.run}, or a workflow-facing service
  operation. Capabilities include side-effect class, cost hint, latency hint,
  reliability hint, and policy labels.
  \item[Agent request.] A request carries a schema version, request id, trace id,
  task id, target capability, serialized input, context entries, timeout budget,
  retry budget, policy labels, and replay mode. The request id is the unit of
  duplicate detection and cooperative cancellation.
  \item[Agent response.] A response carries success/failure, output text,
  structured error code/message, retryability, trace id, and runtime metadata.
  Late responses for cancelled work are rewritten to a cancellation failure.
  \item[State record.] A state record is a namespaced key, kind, scope, value,
  and sequence. The sequence is monotonic within the state store and is used by
  conditional writes for leases and workflow terminal transitions.
  \item[Workflow run.] A workflow run has a run id, workflow identity, source
  name, status, compiled plan, result or error, execution owner, lease record,
  and sequence. Node-level records store per-task progress and outputs.
  \item[Runtime event.] An event has a sequence number, timestamp, trace id,
  component, kind, message, and details. Events record nondeterministic choices,
  model/tool calls, retries, failures, and workflow transitions.
\end{description}

\subsection{State Namespace}

Table~\ref{tab:state-namespace} lists the durable namespaces currently used by
\rasn. The table is important because it shows which decisions are intended to
survive process boundaries and therefore must not be hidden in in-memory helper
objects.

\begin{table}[t]
\caption{Primary \rasn state namespaces.}
\label{tab:state-namespace}
\centering
\begin{tabular}{lll}
\toprule
Key pattern & Kind & Purpose \\
\midrule
\texttt{workflow/<run>} & \texttt{workflow} & Whole-run status and plan \\
\texttt{workflow-node/<run>/<node>} & \texttt{workflow.node} & Node progress/output \\
\texttt{workflow-lease/<run>} & \texttt{workflow.lease} & Execution ownership \\
\texttt{artifact/<id>} & \texttt{artifact.index} & Oversized tool output index \\
\texttt{observability/<id>} & \texttt{observability.snapshot} & Trace snapshot index \\
\texttt{selftest/<id>} & \texttt{selftest} & Self-test state smoke records \\
\bottomrule
\end{tabular}
\end{table}

\subsection{State Record Lifecycle}

State records follow a small lifecycle. They are created by module-specific
writers, validated by the state service, assigned a sequence, optionally written
to the journal, and made visible to reads/queries. Checkpointing snapshots the
current map; recovery loads checkpoint records and replays journal records to
reconstruct the map. Workflow records add a higher-level lifecycle on top of
this generic record lifecycle: \texttt{running} records can be replaced by
terminal records only when the writer still owns the expected sequence, and
lease records can be updated only by the current owner sequence.

This lifecycle is the bridge between an in-process prototype and a future
distributed implementation. The exact persistence backend can change, but the
logical operations are stable: create, compare-and-set update, prefix query,
checkpoint, and recover. The rest of \rasn is written against these operations
rather than against file paths or local maps.

\subsection{Correctness Invariants}

The current prototype enforces the following invariants. They are weaker than a
fully replicated production system, but they are concrete enough to test.

\begin{description}
  \item[I1: Schema discipline.] RPC handlers reject unsupported schema versions
  and missing required identifiers before invoking providers or tools.
  \item[I2: Explicit service ownership.] Service graph startup/shutdown is
  reference counted; one app wrapper cannot stop shared services while another
  wrapper or direct command still owns the graph.
  \item[I3: Capability-routed execution.] Application requests are routed by
  capability descriptors rather than by prompt text or CLI-specific branching in
  the coordinator.
  \item[I4: Side-effect-aware retry.] The coordinator may retry retryable model
  failures within a bounded budget, but it does not retry tool invocations.
  \item[I5: State boundary preservation.] Workflow records, node records, lease
  records, artifact indexes, and observability snapshots are written through
  \texttt{rasn.state}; service mode and direct mode share the same facade.
  \item[I6: Lease-protected workflow execution.] A workflow run must own a
  non-expired execution lease before dispatching nodes. Active owners renew the
  lease; duplicate active execution is rejected; stale ownership can be taken
  over only after expiration.
  \item[I7: Terminal workflow state.] Completed, failed, and cancelled workflow
  states are terminal. A final completion write cannot overwrite a newer
  cancellation record because terminal writes use conditional state sequences.
  \item[I8: Cancellation retention.] Agent cancellation tombstones are retained
  while the corresponding request is still in flight, even when the bounded
  tombstone set is under pressure.
  \item[I9: Error transparency.] State lookup and state-write failures are
  returned as explicit errors, not converted into success-shaped missing-key or
  best-effort results.
  \item[I10: Bounded local effects.] Local tool output is bounded; oversized
  output is moved to artifacts and indexed in state. Write/shell capabilities
  are opt-in.
\end{description}

\subsection{Invariant Enforcement}

Table~\ref{tab:invariant-enforcement} shows how selected invariants are enforced
in the implementation. This is more precise than saying that the system is
``robust'': each row identifies the runtime mechanism and the type of regression
that would violate it.

\begin{table*}[t]
\caption{Selected invariants, enforcement mechanisms, and regression signals.}
\label{tab:invariant-enforcement}
\centering
\begin{tabular}{p{0.18\textwidth}p{0.39\textwidth}p{0.31\textwidth}}
\toprule
Invariant & Enforcement mechanism & Regression signal \\
\midrule
Schema discipline & Version checks and required-field validation in request,
state, workflow, and RPC handlers & Malformed input reaches provider/tool code \\
Side-effect-aware retry & Coordinator retries only retryable model failures and
never retries \texttt{tool.*} & Tool invocation repeated after timeout/failure \\
Lease ownership & Workflow start creates or takes over lease with state CAS;
renewal updates expected lease sequence & Duplicate active run executes nodes \\
Terminal cancellation & Cancel writes terminal state with expected sequence;
completion writes use original running sequence & Cancelled run becomes completed \\
Cancellation retention & Runtime tombstone trimming skips in-flight request ids
& Late cancelled request returns success \\
Error transparency & Recovery and cancel paths propagate non-missing state
errors & State outage reported as missing workflow \\
\bottomrule
\end{tabular}
\end{table*}

\subsection{Request Lifecycle}

The normal model/tool request path follows the lifecycle below. Each step maps
to the service graph, coordinator, registry, or target-agent runtime.

\begin{lstlisting}
1. Application builds typed agent_request.
2. Service graph selects direct call or rDSN clientlet call.
3. Coordinator validates schema and capability.
4. Coordinator queries registry for a healthy descriptor.
5. Coordinator records route decision in the runtime trace.
6. Target agent begins request lifecycle:
   a. reject duplicate request id;
   b. add request id to in-flight set;
   c. check cancellation tombstone;
   d. call provider/tool with timeout budget.
7. Provider/tool returns result or structured failure.
8. Runtime converts late cancelled result to request_cancelled.
9. Coordinator records retry/failure/finish events and returns response.
\end{lstlisting}

The lifecycle deliberately separates routing, execution, and cancellation. This
is the distributed-systems lesson carried into \rasn: each boundary has a clear
owner and a clear failure mode, which makes it possible to debug and harden one
boundary at a time.

\subsection{Timeout, Retry, and Cancellation Semantics}

Timeouts, retries, and cancellation are related but not equivalent. A timeout is
a budget attached to a request or workflow node; in service mode it becomes an
\rdsn client-side deadline and in provider mode it is propagated to provider
configuration where possible. A retry is a coordinator decision after a
retryable model failure; it consumes retry budget and records a retry event. A
cancellation is a caller decision that marks a request or workflow run as no
longer desired; it must prevent a later successful result from being surfaced.

\rasn deliberately avoids conflating these mechanisms. A timed-out tool is not
retried because it may already have performed a side effect. A cancelled
provider call is not forcibly killed by the current prototype, but the runtime
records the cancellation and rewrites the eventual response. A workflow cancel
is durable: it writes the \texttt{cancelled} state through \texttt{rasn.state}
and therefore survives process restart. This separation gives the runtime
well-defined behavior even when the underlying provider or shell command is not
preemptible.

\subsection{Replay and Diagnosis Semantics}

Replay in the current system is best understood as incremental deterministic
re-execution plus trace-guided diagnosis rather than full process replay. A
replay load can read prior trace events and expose nondeterministic choices,
provider calls, workflow transitions, state decisions, and tool outcomes. The
runtime can reuse simulator choices, recorded \texttt{llm.response} events in
provider order, and matching recorded \texttt{tool.ok} or \texttt{tool.error}
events. Workflow execution also records \texttt{workflow.node.start} and
\texttt{workflow.node.finish} transitions. During replay, the executor checks
the recorded node-start order before dispatching a node; a mismatch is reported
as a \texttt{replay.miss} failure before any model or tool side effect occurs.
CodePilot file tools record \texttt{filesystem.snapshot} events for
read/list/search targets. Replay checks those path fingerprints before a file
tool consumes current workspace state, surfacing filesystem drift rather than
silently reusing stale file context.
When replay mode is active and a side-effecting tool has no matching recorded
result, \texttt{rasn.tool.agent} fails closed instead of repeating a write or
shell command. Side-effecting tool intents also record
\texttt{external.effect} events that carry effect class, operation, stable
fingerprint, replay policy, and status for commits, policy denials, replay hits,
and replay misses. In-process timing and randomness already draw from rDSN's
pluggable environment provider, so the remaining virtualization gap is the
external environment. Full replay would still require stronger virtualization for
external environments, but the current model now gives non-file effects an
auditable runtime ledger instead of only a transcript. The current model
therefore advances replay for concrete model, tool, scheduler, local-file, and
external-effect-intent decisions while leaving broader environment virtualization
as a future hardening target.

\subsection{Workflow State Machine}

Workflow execution is the most complete example of the runtime model because it
combines compilation, routing, state, leases, and observability. A run moves
through the states in Figure~\ref{fig:workflow-state-machine}.

\begin{figure}[t]
\centering
\begin{lstlisting}
new
  | validate + compile + acquire lease
  v
running
  | all nodes ok         | node/tool/model error
  v                      v
completed              failed
  ^
  |
cancel requested or observed between nodes
  |
cancelled
\end{lstlisting}
\caption{Workflow run state machine. Terminal states are \texttt{completed},
\texttt{failed}, and \texttt{cancelled}; terminal transitions are persisted
through \texttt{rasn.state}.}
\Description{A text diagram showing workflow states new, running, completed,
failed, and cancelled, with validation and lease acquisition before running and
terminal transitions persisted through state.}
\label{fig:workflow-state-machine}
\end{figure}

The implementation does not assume that the workflow process that receives a
cancel or query request is the same process that originally started the run.
Both operations can recover \texttt{workflow/<run-id>} from state. This makes
workflow management restart-safe even before the state store itself is replaced
by a replicated backend.

Each transition in Figure~\ref{fig:workflow-state-machine} has a concrete
guard:

\begin{description}
  \item[\texttt{new} to \texttt{running}.] The workflow source parses, the graph
  is acyclic, resume/source identity checks pass, and the service acquires the
  execution lease.
  \item[\texttt{running} to \texttt{completed}.] All executable nodes complete,
  no cancellation is observed, and the terminal write matches the running
  record's expected sequence.
  \item[\texttt{running} to \texttt{failed}.] A node fails or a required runtime
  boundary returns a non-retryable error; dependent nodes become blocked.
  \item[\texttt{running} to \texttt{cancelled}.] A cancel request writes a
  terminal record, or execution observes the cancelled record between nodes.
  \item[Restart recovery.] Query, cancel, and resume may recover the run record
  from state when it is absent from memory. Non-missing state errors abort the
  operation.
\end{description}

The workflow state machine is therefore not merely conceptual. It determines
which state writes require compare-and-set, which errors are safe to retry, and
which CLI commands can be implemented after a restart.

\section{Architecture}

Figure~\ref{fig:architecture} summarizes the current architecture.

\begin{figure}[t]
\centering
\begin{lstlisting}
CodePilot / future applications
        |
        v
rasn_service_graph
        |
        +--> rasn.coordinator <----> rasn.registry
        |          |
        |          +--> rasn.llm.agent  --> LLM providers
        |          +--> rasn.tool.agent --> tool providers/policy
        |
        +--> rasn.workflow ----> rasn.state
        +--> rasn.observability -> rasn.state snapshots
\end{lstlisting}
\caption{High-level \rasn service graph. Direct mode calls these components
in-process; service mode routes through typed \rdsn RPC task codes.}
\Description{A layered architecture where CodePilot and future applications
call a rASN service graph, which routes to coordinator, registry, model agent,
tool agent, workflow service, state service, and observability service.}
\label{fig:architecture}
\end{figure}

\subsection{Service Graph and Lifecycle}

\texttt{rasn\_service\_graph} is the facade used by applications. It owns the
in-process service implementations and the addresses of service-mode RPC
endpoints. It provides operations such as \texttt{complete},
\texttt{run\_tool}, \texttt{invoke}, \texttt{put\_state},
\texttt{start\_workflow}, and \texttt{query\_events}.

The graph has reference-counted lifecycle ownership. Each \rdsn app wrapper
that depends on it acquires a lifecycle reference during startup and releases
it during shutdown. Direct \codepilot commands acquire a command-scoped
reference. This prevents one app wrapper from stopping shared services while
other services are still active. Startup and shutdown work is performed outside
the lifecycle lock to avoid lock-held RPC registration and deadlock-prone waits.

\subsection{RPC Surface}

\rasn defines explicit \rdsn task codes for generic agent control and for each
service family:

\begin{itemize}
  \item Generic agent-control RPCs for describe, invoke, cancel, heartbeat, and
  query.
  \item Registry RPCs for register, unregister, heartbeat, query, and list.
  \item State RPCs for put, conditional put, get, query, checkpoint, and
  recovery.
  \item Workflow RPCs for validate, compile, start, query, and cancel.
  \item Model gateway RPCs for provider description, provider selection, and
  health.
  \item Observability RPCs for events, failures, replay loading, and snapshots.
\end{itemize}

Service-mode \codepilot waits for all required service surfaces to respond
before dispatching a command. The readiness gate probes state, registry,
coordinator, model-agent identity, model health, tool-agent identity, workflow
validation, and observability. Failures name the specific service boundary.

\subsection{Execution Modes}

\rasn supports two execution modes that intentionally share the same logical
module interfaces.

\begin{description}
  \item[Direct mode.] \codepilot or a test owns a command-scoped
  \texttt{rasn\_service\_graph} reference. Service calls are in-process C++
  calls, but they still pass through the same graph facade, policy checks,
  runtime validation, state writers, and observability hooks.
  \item[Service mode.] \rdsn starts service apps for registry, coordinator,
  model agent, tool agent, state, workflow, observability, and CodePilot. Calls
  cross typed \rdsn RPCs using clientlets and serverlets. This mode exercises
  serialization, timeout propagation, startup ordering, readiness probes, and
  per-service failure visibility.
\end{description}

The design goal is not to maintain two implementations. Direct mode is a fast
developer path; service mode is the distributed-systems path. Bugs found in one
mode should usually be fixed in the shared service graph or module store rather
than by adding mode-specific branches.

\subsection{Module Contracts}

Table~\ref{tab:module-contracts} summarizes each module's contract. The most
important column is state ownership: it identifies which runtime facts the
module is allowed to mutate directly and which facts must cross a service
boundary.

\begin{table*}[t]
\caption{\rasn module contracts.}
\label{tab:module-contracts}
\centering
\small
\begin{tabular}{p{0.13\textwidth}p{0.19\textwidth}p{0.25\textwidth}p{0.29\textwidth}}
\toprule
Module & Inputs & Owned state & Outputs/failures \\
\midrule
Runtime & Agent requests, config & Lifecycle, in-flight ids, tombstones & Responses, trace events \\
Registry & Descriptors, heartbeats & Descriptor table, lease timestamps & Capability matches, stale lease removal \\
Coordinator & Capability requests & Route decisions, retry counters & Agent responses, retry/failure events \\
Model agent & Completion requests & Provider selection, health & Text responses, provider errors \\
Tool agent & Tool requests & Tool provider config & Tool output, policy/tool errors \\
Policy & Tool metadata, config & Artifact indexes through state & Allow/deny, spill metadata \\
State & State put/get/query & Records, sequence, checkpoint, journal & CAS success/conflict, recovery errors \\
Workflow & Workflow source/query & Runs, nodes, leases through state & Plans, run records, terminal status \\
Observability & Runtime events, queries & Event buffer, snapshot indexes & Events, failures, diagnostics \\
CodePilot & CLI commands/config & Session-local command context & User output, structured errors \\
\bottomrule
\end{tabular}
\end{table*}

\subsection{Control Flow}

The following control-flow sketches describe the two most common paths. First,
an interactive CodePilot prompt becomes a model request:

\begin{lstlisting}
codepilot prompt
  -> build agent_completion_request
  -> rasn_service_graph.complete()
  -> coordinator.invoke(capability=model.complete)
  -> registry.query(model.complete)
  -> model agent invoke
  -> selected llm_provider.complete()
  -> response + trace events
\end{lstlisting}

Second, a declarative workflow run becomes a sequence of state-backed node
executions:

\begin{lstlisting}
workflow start file
  -> validate/compile workflow_graph
  -> acquire workflow-lease/<run-id> by conditional put
  -> persist workflow/<run-id> = running
  -> for each ready node:
       persist workflow-node/<run>/<node> = running
       dispatch model/tool through coordinator
       persist workflow-node/<run>/<node> = completed/failed
  -> conditional terminal write to workflow/<run-id>
  -> release lease
\end{lstlisting}

These flows expose where failures can be observed: registry lookup failures,
provider failures, policy denials, state conflicts, lease ownership failures,
and terminal state persistence failures are separate error classes rather than a
single CLI failure string.

\subsection{Module Decomposition and Ownership}

\rasn is decomposed so that each module owns one class of runtime state and
exposes that state through typed operations rather than shared global variables
or prompt conventions.

\begin{description}
  \item[Task and message model.] The schema layer defines versioned requests,
  responses, errors, descriptors, metadata entries, model messages, tool
  requests, workflow records, state records, and observability events. This
  layer is intentionally passive: it validates shape and serializes through
  \rdsn binary readers/writers, but it does not make routing or policy
  decisions.
  \item[Agent runtime.] The runtime owns lifecycle state, identity, capability
  descriptors, request validation, trace emission, deadline defaults, and
  cooperative cancellation tracking for service-like agents.
  \item[Registry.] The registry owns the mapping from capabilities to live
  descriptors. It separates static configured agents from dynamic agents that
  must refresh heartbeat leases.
  \item[Coordinator.] The coordinator owns capability routing, retry decisions,
  side-effect-aware dispatch, and route observability. It is the point where
  application-level requests become agent invocations.
  \item[Model agent.] The model agent isolates LLM provider configuration and
  provider health from the rest of the runtime. It exposes model completion as a
  normal capability, whether the backing provider is a simulator, cloud API, or
  local HTTP-compatible server.
  \item[Tool agent and policy.] The tool agent owns local side-effect execution.
  The policy manager classifies operations, blocks unsafe defaults, bounds
  output, and records oversized artifact metadata through the state service.
  \item[State service.] The state service owns durable runtime records:
  checkpoints, journals, workflow run state, node state, leases, artifact
  indexes, and observability snapshot indexes.
  \item[Workflow service.] The workflow service owns declarative workflow
  validation, graph compilation, execution leases, per-node progress,
  cancellation, and resume.
  \item[Observability.] The observability module owns runtime event capture,
  failure summaries, replay metadata, snapshots, and diagnostic queries.
  \item[Applications.] \codepilot and future applications own user interface,
  prompts, domain-specific skills, and command syntax. They call \rasn services
  instead of embedding runtime state machines in the CLI.
\end{description}

\subsection{Registry and Routing}

The registry stores agent descriptors containing agent ID, role, app name,
host/port or URI endpoint, version, health, and capabilities. Static external
agents can be configured with \texttt{[rasn.agent.*]} sections. Built-in agents
can dynamically register and send periodic heartbeats in service mode. Registry
leases are swept by a timer so stale dynamic descriptors are removed while
static descriptors remain available.

The coordinator resolves a request capability to an agent descriptor, records
the route decision, and invokes the selected agent. Model-agent failures can be
retried when they are marked retryable and the request carries a bounded retry
budget. Tool capabilities are not retried because a timed-out tool invocation
may already have performed side effects.

\subsection{State, Checkpoints, and Leases}

The state service provides a small key-value store with namespaced keys,
sequence numbers, checkpoints, append-only journals, and conditional put. It is
not intended as a production replicated database in the current prototype.
Instead, it defines the boundary through which runtime state flows. Workflow
execution uses conditional writes to acquire and update
\texttt{workflow-lease/<run-id>} records. Active workflow owners renew leases
while nodes execute. Stale leases can be taken over after the configured lease
TTL. The workflow service now distinguishes missing state keys from state
lookup failures during run recovery.

\subsection{Workflow Runtime}

The workflow subsystem parses a lightweight declarative workflow language into
a graph. Nodes represent model or tool tasks with dependencies, policy labels,
execution budgets, retry budgets, latency hints, cost hints, reliability hints,
and optional state keys. The compiler validates dependencies and rejects cycles.
An optimizer orders ready nodes by score and reports stages, maximum
parallelism, critical-path latency, cost, and minimum reliability. Execution
persists whole-run and per-node state. Resume mode loads completed node outputs
from state and executes incomplete downstream nodes.

Workflow cancellation is cooperative. A cancel request marks the run cancelled,
remaining nodes are persisted as cancelled, and final completion cannot
overwrite a cancelled run. Agent-level cancellation is also cooperative:
currently it tracks in-flight request IDs and converts late results into
\texttt{request\_cancelled}; it does not yet forcibly terminate arbitrary
provider or shell work.

\subsection{Observability}

\rasn's runtime log records events as JSON lines when tracing is configured.
The observability service can query events, failures, timelines, diagnostic
summaries, replay loads, and snapshots. Snapshot metadata is persisted through
the state service, giving operators a durable index of trace file, event count,
failure count, and last sequence.

Beyond per-run traces, \rasn aggregates quantitative runtime metrics through
\rdsn's existing \texttt{perf\_counter} subsystem rather than a parallel metrics
stack. A process-global registry maintains cumulative counters per runtime event
kind, lazily created per-failure-class counters, and task, model, and tool
latency percentile counters, all in the \texttt{rasn} counter section. The same
snapshot is rendered as an aligned text table, Prometheus exposition format, or
JSON, and is reachable both through the \codepilot \texttt{observe metrics}
command and through a \texttt{rasn.metrics} command registered with \rdsn's
\texttt{command\_manager} so a running deployment answers over the local or
remote \rdsn CLI.

Metrics make a degrading dependency visible; a circuit breaker makes the runtime
react to it. The model gateway wraps each provider in a per-provider
consecutive-failure circuit breaker (\texttt{circuit\_breaker.*}): after a
configurable number of consecutive failures it opens and short-circuits requests
for a cooldown, then admits a single half-open probe that either closes or
reopens it. The breaker fast-fails a call to a dead or hanging endpoint
\emph{before} the provider's own retry loop, so latency and budget are not wasted
on it. It draws its clock from \texttt{dsn\_now\_ms} (the same pluggable
\texttt{env\_provider} used elsewhere), so breaker timing is deterministic under
replay and replayed runs and in-process providers (the simulator and other
providers that perform no network I/O) bypass it entirely, while network-backed
providers are guarded even when their endpoint is loopback (Ollama, llama.cpp, LM
Studio).
Breaker activity flows through the same \texttt{record\_event} choke point as
every other metric, and live per-provider state is reachable through \codepilot's
\texttt{observe resilience} and a \texttt{rasn.resilience} \texttt{command\_manager}
command.

A circuit breaker protects against a \emph{broken} dependency; admission control
protects against overwhelming a \emph{healthy} one. Each provider is therefore also
fronted by an admission gate (\texttt{admission\_gate.*}): a concurrency bulkhead
that caps simultaneous in-flight requests (excess load fast-fails) plus graceful
backpressure that, once load passes a soft threshold, applies a short pre-call delay
bounded by a configured ceiling. The backpressure curve reuses \rdsn's own
\texttt{exp\_delay} admission primitive --- the staged-delay mechanism the host
runtime uses to throttle its task queues --- rather than a hand-rolled schedule. The
gate is evaluated before the breaker and after the replay check, and shares the
breaker's \texttt{in\_process()} exemption, so replayed runs, the simulator, and
sequential single-agent use are unaffected; it engages only under the concurrent
fan-out of service mode. Rejections and backpressure delays are reported through the
same two-counter pattern as the breaker, and live in-flight depth appears in the
\texttt{rasn.resilience} command.

A breaker covers failure and admission control covers concurrency; the third
dimension, throughput, is covered by a per-provider token-bucket rate limiter
(\texttt{rate\_limiter.*}). Hosted model APIs enforce requests-per-minute quotas,
and a concurrency bulkhead does not bound them --- one sequential caller stays at
in-flight $\approx 1$ yet can still exceed a per-minute limit and trigger HTTP~429s
that waste the retry budget. The limiter paces calls under the configured rate,
delaying a request briefly (reserving the next token) when slightly ahead and
fast-failing when the wait would exceed a bound. Like the breaker it draws its
refill clock from \texttt{dsn\_now\_ms}, so it is replay-deterministic; it sits
beside the other two gates (acquiring before the breaker so a rate rejection never
strands a half-open probe), shares the \texttt{in\_process()} exemption, and is
disabled by default so existing behavior is unchanged. Its two counters follow the
same pattern, and live per-provider rate, burst, and tokens appear in the
\texttt{rasn.resilience} command.

The tool gateway reuses the same admission and rate engines, keyed by tool name.
After replay lookup and policy approval, \texttt{rasn.tool.agent} reserves an
admission slot, optionally acquires a rate token, coalesces the requested delays,
and then invokes the registered provider. Replayed results, replay-missing
side-effect failures, and policy denials bypass the gates, so robustness controls do
not perturb replay or safety semantics. The service now holds its provider
registration lock only while taking a shared provider reference; actual tool
execution is governed by explicit per-tool limits rather than an implicit global
mutex. Four additional counters
(\texttt{rasn\_tool\_admission\_rejected\_total},
\texttt{rasn\_tool\_admission\_delayed\_total},
\texttt{rasn\_tool\_rate\_limited\_total}, and
\texttt{rasn\_tool\_rate\_delayed\_total}) and live per-tool state are exposed
through the same \texttt{rasn.resilience} surface.

Finally, service-mode coordinator dispatch is treated as its own remote
dependency boundary. After the registry selects a remote agent for a capability,
\texttt{rasn\_coordinator\_service} now guards the outbound
\texttt{RPC\_RASN\_AGENT\_INVOKE} path per remote \texttt{agent\_id} with the same
breaker/admission/rate trio. Only retryable dependency outcomes --- transport
errors and retryable remote responses --- count against the breaker; deterministic
application responses such as policy denials, validation failures, or tool errors
do not poison the agent. The guards run only on the RPC-client path, so inline CLI
behavior is unchanged. Seven additional counters
(\texttt{rasn\_remote\_agent\_breaker\_open\_total},
\texttt{rasn\_remote\_agent\_breaker\_short\_circuit\_total},
\texttt{rasn\_remote\_agent\_admission\_rejected\_total},
\texttt{rasn\_remote\_agent\_admission\_delayed\_total},
\texttt{rasn\_remote\_agent\_rate\_limited\_total},
\texttt{rasn\_remote\_agent\_rate\_delayed\_total}, and
\texttt{rasn\_remote\_agent\_endpoint\_invalid\_total}) and live per-agent state are
reported through \texttt{observe resilience} and \texttt{rasn.resilience}.

\section{Implementation Details}

\subsection{Code Organization}

The \rasn source tree is organized as a \rdsn plugin:

\begin{lstlisting}
src/plugins/rasn/
  agent_types.h
  agent_messages.h
  agent_runtime.{h,cpp}
  agent_registry.{h,cpp}
  agent_clients.{h,cpp}
  coordinator_service.{h,cpp}
  agent_services.{h,cpp}
  state_service.{h,cpp}
  workflow.{h,cpp}
  workflow_service.{h,cpp}
  policy_manager.{h,cpp}
  observability.{h,cpp}
  metrics.{h,cpp}
  circuit_breaker.{h,cpp}
  admission_gate.{h,cpp}
  rate_limiter.{h,cpp}
  llm_provider.{h,cpp}
  agent_tools.{h,cpp}
  codepilot/
  tests/
  examples/
  docs/report/
\end{lstlisting}

The build is organized so that rASN is a reusable platform rather than a
single application binary. The engine sources in \texttt{src/plugins/rasn/}
compile into a \texttt{rasn} static library, and both the \codepilot
executable (\texttt{codepilot/}, including \texttt{codepilot/main.cpp}) and the
\texttt{rasn.unit\_tests} executable link that library instead of recompiling
the engine. \texttt{CMakeLists.txt} integrates the library, the application,
and the unit tests into the \rdsn build. The implementation is intentionally
header-heavy for message serialization helpers and uses \cpp service classes
for runtime behavior.

Table~\ref{tab:file-responsibilities} maps the main implementation files to
their responsibilities. This table is also useful for future maintainers because
it distinguishes reusable nucleus code from the CodePilot application layer.

\begin{table*}[t]
\caption{Main source files and responsibilities.}
\label{tab:file-responsibilities}
\centering
\begin{tabular}{lll}
\toprule
Files & Layer & Responsibility \\
\midrule
\texttt{agent\_types/messages} & Core schema & Versioned descriptors, requests, responses, serialization \\
\texttt{agent\_runtime} & Core runtime & Agent lifecycle, validation, traces, cancellation \\
\texttt{agent\_registry} & Service & Descriptor table, heartbeat leases, sweeps \\
\texttt{coordinator\_service} & Service & Capability routing, bounded retry, dispatch \\
\texttt{agent\_clients} & RPC client & Typed clientlet wrappers and deadlines \\
\texttt{agent\_services} & Composition & Service graph, app wrappers, direct/RPC switching \\
\texttt{state\_service} & Service & Key-value state, CAS, checkpoints, journals, recovery \\
\texttt{workflow/workflow\_service} & Service & Workflow parser, optimizer, executor, leases, resume/cancel \\
\texttt{policy\_manager/redaction} & Policy & Tool classification, gating, redaction, artifact indexing \\
\texttt{observability} & Service & Events, failures, snapshots, replay metadata \\
\texttt{metrics} & Observability & Perf-counter metric registry and text/Prometheus/JSON rendering \\
\texttt{circuit\_breaker} & Resilience & Per-provider model circuit breaker (closed/open/half-open) \\
\texttt{admission\_gate} & Resilience & Per-provider/per-tool concurrency bulkhead and \texttt{exp\_delay} backpressure \\
\texttt{rate\_limiter} & Resilience & Per-provider/per-tool token-bucket rate limiter (requests/minute pacing) \\
\texttt{llm\_provider/model\_agent} & Provider & Simulator and HTTP-compatible model adapters \\
\texttt{agent\_tools} & Provider API & Generic tool provider interface \\
\texttt{codepilot/*} & Application & CLI commands, local tools, skills, provider configuration \\
\bottomrule
\end{tabular}
\end{table*}

\subsection{rDSN Integration Points}

\rasn intentionally reuses rDSN mechanisms instead of implementing a parallel
runtime. The important integration points are:

\begin{itemize}
  \item \texttt{dsn::service\_app} for long-lived roles such as
  \texttt{rasn.workflow}, \texttt{rasn.state}, and \texttt{rasn.codepilot}.
  \item \texttt{dsn::serverlet<T>} and \texttt{dsn::rpc\_replier<T>} for typed
  RPC handlers.
  \item \texttt{dsn::clientlet} wrappers for synchronous service-mode calls with
  propagated timeout budgets.
  \item \texttt{DEFINE\_TASK\_CODE\_RPC} task codes for every service boundary
  and control operation.
  \item \texttt{binary\_writer} and \texttt{binary\_reader} for message
  serialization.
  \item \texttt{dsn\_config\_get\_value\_*} for provider, policy, lease,
  timeout, endpoint, and persistence settings.
  \item \texttt{dsn\_now\_ms}/\texttt{dsn\_random64} (the pluggable
  \texttt{env\_provider} for time and randomness) for runtime timing, latency
  measurement, and identifier/sampling randomness, so these nondeterministic
  inputs are virtualizable and seedable for deterministic replay rather than
  drawn from private wall-clock-seeded generators.
  \item \texttt{dsn::service::zlock} for shared state in stores and runtimes.
  \item \texttt{dsn::tasking::enqueue\_timer} for registry heartbeats, lease
  sweeps, delayed workflow recovery, and workflow lease renewal.
  \item \texttt{dsn::utils::filesystem} for checkpoints, temporary files,
  artifact paths, and cross-platform file operations.
  \item \texttt{perf\_counter} (the C \texttt{dsn\_perf\_counter\_*} API backed by
  \texttt{simple\_perf\_counter}) for cumulative runtime counters and task,
  model, and tool latency percentiles in the \texttt{rasn} counter section.
  \item \texttt{command\_manager} (\texttt{register\_command}/\texttt{run\_command})
  to expose the \texttt{rasn.metrics} and \texttt{rasn.resilience} operational
  commands over the \rdsn local and remote CLI.
  \item \texttt{exp\_delay} (\texttt{include/dsn/utility/exp\_delay.h}), \rdsn's
  staged admission-control delay, reused as the model and tool gateways'
  backpressure curve so overload throttling matches the host runtime's own
  task-queue policy.
  \item \texttt{dsn\_now\_ms} (the env-provider clock) injected into the model
  circuit breaker and model/tool token-bucket rate limiters, so cooldown and token
  refill are deterministic under replay and seedable rather than reading a private
  wall clock.
\end{itemize}

The implementation currently still has same-process globals for convenience,
but service-mode calls deliberately cross \rdsn client/server boundaries where
possible. This is a stepping stone toward deployment as independently hosted
services.

\subsection{Task and Message Schema}

\texttt{agent\_types.h} and \texttt{agent\_messages.h} define the common data
model shared by the runtime, service wrappers, and applications. Important
records include agent descriptors, capabilities, requests, responses, errors,
completion requests, tool requests, trace metadata, and key-value metadata
entries. Each top-level structure carries \texttt{schema\_version}; validation
rejects unsupported versions and missing required identifiers before a request
can reach model or tool code.

The schema is deliberately explicit. A model call is not just a prompt string:
it carries a task id, trace id, request id, capability name, policy labels,
timeout budget, retry budget, context entries, and replay mode. Likewise, a tool
call is a typed request with a tool name and argument vector rather than shell
text hidden inside a prompt. The implementation provides \rdsn
\texttt{marshall}/\texttt{unmarshall} functions for these records, allowing the
same messages to be used in direct in-process mode and in service-mode RPC.

Table~\ref{tab:request-fields} shows why these fields are part of the runtime
protocol rather than application-level annotations.

\begin{table}[t]
\caption{Selected request fields and their runtime role.}
\label{tab:request-fields}
\centering
\begin{tabular}{ll}
\toprule
Field & Runtime role \\
\midrule
\texttt{request\_id} & Duplicate detection and cancellation key \\
\texttt{trace\_id} & Correlates model, tool, state, and workflow events \\
\texttt{capability} & Registry lookup and coordinator routing key \\
\texttt{timeout\_ms} & Client RPC deadline and provider budget \\
\texttt{retry\_budget} & Bounded model retry control \\
\texttt{policy\_labels} & Tool/model policy context \\
\texttt{replay\_mode} & Hook for trace-guided replay \\
\bottomrule
\end{tabular}
\end{table}

Validation is intentionally performed before side effects. For example, an
unsupported schema version, empty capability, missing request id, or malformed
tool request is rejected by the runtime/coordinator rather than reaching a
provider. This rule is what lets tests assert deterministic failure strings for
bad inputs and prevents providers from becoming implicit validators.

\texttt{schema\_manifest.*} now derives multiple integration artifacts from the
same descriptor table. The text output is documentation, the JSON output is
suited for dashboards or code generators, the IDL output is a compact
language-neutral contract, the C++ output is a generated SDK header, the C++ RPC
client output is a generated \rdsn \texttt{clientlet} layer, the TypeScript
output is a set of exported interfaces, the TypeScript RPC-client output is a
transport-backed client class layer, the Python output is a dataclass module,
and the Python RPC-client output is a transport-backed client class layer. The
C++ data-contract generator maps rASN scalar types to
\texttt{std::string}, \texttt{bool}, \texttt{std::uint32\_t}, and
\texttt{std::uint64\_t}; the TypeScript generator maps them to
\texttt{string}, \texttt{boolean}, and \texttt{number}; and the Python generator
maps them to \texttt{str}, \texttt{bool}, and \texttt{int}. Repeated fields are
mapped to \texttt{std::vector<T>}, \texttt{T[]}, or \texttt{List[T]}; Python
uses \texttt{field(default\_factory=list)} so generated dataclasses do not share
mutable defaults. Because the generators are manifest-driven, adding a field to
a runtime contract also changes all exported formats and the regression test can
detect missing nested records before an external integration relies on an
incomplete SDK surface. The RPC-client generators are operation-manifest driven:
C++ methods call the matching \texttt{RPC\_RASN\_*} task code, preserve
timeout/thread/partition parameters, and return
\texttt{std::pair<error\_code,T>} like the hand-written clients; TypeScript and
Python methods call a host-supplied transport with the same task-code constant
and typed request/response pair. The
executable handles \texttt{schema} as a direct fast path before starting the
\rdsn runtime, so redirected schema artifacts are not polluted by service
startup or shutdown logs.

\subsection{Service Graph and RPC Wiring}

\texttt{agent\_services.*} contains the shared service graph and the rDSN
service-app wrappers. \texttt{rasn\_service\_graph} is the facade used by
\codepilot and tests. In direct mode it calls in-process service objects; in
service mode it creates typed \rdsn clientlets for registry, coordinator, model,
tool, state, workflow, and observability endpoints. The graph also installs
state writers for policy and observability modules so artifact and snapshot
indexes cross the same state boundary in both modes.

Service wrappers use \texttt{dsn::service\_app} for lifecycle and
\texttt{dsn::serverlet<T>} for RPC handlers. Each handler is intentionally thin:
it decodes a typed request, calls the corresponding store/runtime method, and
replies with a typed response. This keeps most correctness logic outside the RPC
boilerplate while still exercising rDSN task codes, serialization, thread pools,
and client timeouts.

The service graph enforces startup discipline through readiness checks. The
current readiness algorithm is:

\begin{lstlisting}
for each required service boundary:
  send a cheap typed probe with a short deadline
  if probe fails:
    return "service <name> not ready: <cause>"
dispatch user command only after all probes pass
\end{lstlisting}

The probes are intentionally semantic rather than just socket-level checks:
state must answer a prefix query, registry must list/query, the model agent must
describe itself and report health, workflow must validate an empty-safe probe,
and observability must answer a query. This catches miswired service graphs that
would pass a port-open test but fail at the first real command.

\subsection{Agent Runtime}

\texttt{agent\_runtime} is the base class for built-in service agents. It owns
the descriptor, capabilities, lifecycle state, validation, rejection, default
timeouts, trace emission, and request cancellation tracking. Capabilities carry
task type, side-effect class, cost hint, latency hint, reliability hint, and
policy labels. The runtime uses these descriptors to answer describe/query RPCs
and to make registry entries meaningful to the coordinator.

The cancellation tracker maintains an in-flight request set and a bounded set of
cancellation tombstones. A duplicate request id is rejected while work is active;
a cancellation marks the id; and a provider result that arrives after
cancellation is converted into a structured \texttt{request\_cancelled}
response. Tombstone trimming skips still-in-flight requests so cancellation
remains terminal under cancel storms. This invariant is covered by a unit test
that keeps one request active while forcing more than 1024 other cancellations.

The runtime's local request state machine is:

\begin{lstlisting}
absent
  -> begin_request(request_id)
in_flight
  -> cancel_request(request_id)
cancelled_in_flight
  -> finish_request(request_id)
cancelled_tombstone
  -> trim only after request is no longer in flight
\end{lstlisting}

This state machine does not pretend to preempt arbitrary provider code. Instead,
it provides a reliable semantic contract at the agent boundary: once the runtime
acknowledges cancellation for a known in-flight id, the caller will not later see
that request as successful.

\subsection{Registry and Coordinator}

\texttt{agent\_registry.*} implements descriptor registration, unregister,
heartbeat, list, and capability query. Static descriptors come from
configuration; dynamic descriptors are refreshed by heartbeat RPCs and include a
lease timestamp. A sweep timer removes expired dynamic entries but leaves static
entries intact. This mirrors distributed-service membership: a configured
external service remains addressable, while a service discovered through
heartbeat must prove liveness.

\texttt{coordinator\_service.*} is the routing boundary. It receives generic
agent requests, resolves the requested capability through the registry, records
the route decision in the runtime trace, and invokes the selected agent through
the service graph. Retry is intentionally conservative. Only retryable model
failures are retried, retry count is capped by configuration, and tool
capabilities are not retried because a timed-out tool may already have changed
local state.

Registry query results are filtered by role, health, capability, and lease
freshness. The current prototype uses simple first-match routing after filtering;
future versions can replace this with cost/latency/reliability-aware placement
without changing the request schema. Coordinator retry follows this policy:

\begin{lstlisting}
attempt = 0
while true:
  route = registry.query(capability)
  response = invoke(route, request)
  if response.ok: return response
  if capability is tool.*: return response
  if not response.retryable: return response
  if attempt >= min(request.retry_budget, config.max_retry_budget):
    return response
  record retry event
  attempt += 1
\end{lstlisting}

Separating registry filtering from retry policy matters for diagnosability. A
failure can be attributed to no matching capability, an unhealthy/stale agent, an
RPC timeout, a provider-level retryable error, or a non-retryable tool/policy
failure.

\subsection{Model Agent and Provider Adapters}

The LLM provider interface is:

\begin{lstlisting}[language=C++]
class llm_provider {
public:
  virtual std::string name() const = 0;
  virtual std::string model() const = 0;
  virtual model_provider_descriptor describe() const = 0;
  virtual llm_response complete(const llm_request&,
                                nucleus_runtime&) = 0;
  virtual llm_response complete_streaming(
      const llm_request&, nucleus_runtime&,
      const llm_stream_callback& on_chunk);
};
\end{lstlisting}

The prototype includes a local simulator and curl-based HTTP providers for
OpenAI-compatible, Copilot-compatible, Ollama, llama.cpp, and LM Studio style
endpoints. Request budgets are propagated into curl \texttt{max-time} and
capped by provider configuration. Tokens now flow through explicit
\texttt{token\_ref} credential handles: \texttt{env:} references environment
variables, \texttt{file:} reads a local secret file, and \texttt{cmd:} invokes a
command at request time. Legacy \texttt{token\_env} and
\texttt{token\_command} settings remain as compatibility paths. Provider
descriptors expose handles rather than token values or command output, and
command-backed handles are summarized as \texttt{cmd:<configured>}. Production
deployments should still replace the lightweight curl layer with provider SDKs
or a hardened HTTP client.

\texttt{model\_agent.*} adapts the provider interface to the generic agent RPC
surface. Provider selection and health are exposed as explicit model-gateway
operations rather than hidden in CLI configuration. The simulator provider is
important for testing: it can return random-looking local output without network
access, and its choices are recorded as nondeterministic runtime events so replay
experiments have a concrete hook.

The provider interface now includes a provider-neutral streaming method. Native
streaming providers can call the supplied callback as chunks arrive. Providers
without native streaming inherit a default implementation that chunks the final
completion, which gives applications a stable streaming API even when the
underlying provider does not yet expose incremental delivery. Each emitted chunk
is redacted and recorded as an \texttt{llm.response.chunk} event before the
application sees it. This makes streaming observable and replay-diagnosable
instead of a UI-only side channel.

The HTTP provider layer currently uses curl as a portability-first adapter. It
constructs provider-specific request bodies, applies configured headers, writes
temporary response files, bounds request duration, and parses a small set of
common response shapes. Supported deployment styles include:

\begin{description}
  \item[Copilot-compatible.] A configured endpoint and token source compatible
  with GitHub Copilot-style chat completion APIs.
  \item[OpenAI-compatible.] Generic chat-completion endpoints used by many cloud
  and proxy providers.
  \item[Ollama.] Local HTTP serving with a model name and localhost endpoint.
  \item[llama.cpp.] Local server mode for self-hosted model experiments.
  \item[LM Studio.] Local OpenAI-compatible server mode for desktop model
  serving.
  \item[Simulator.] Offline provider that avoids network dependencies and
  records nondeterministic choices.
\end{description}

Provider health is represented as runtime data, not just logs. CodePilot can ask
the model gateway for provider name, model, locality, and health before issuing
an expensive request. This makes provider misconfiguration a readiness failure
instead of a mysterious model-call failure later in a workflow.

\subsection{Tool Agent and Policy Manager}

\texttt{agent\_tools.*} and \texttt{codepilot/local\_tools.*} implement local
tool execution behind the generic \texttt{tool.run} capability. The current tool
provider supports listing, reading, searching, writing, replacing, and shell
execution. The \texttt{policy\_manager} classifies each operation by capability
and side-effect class. Read operations require existing targets; write, replace,
shell, and network-like operations are disabled unless configuration enables
them. Phase 48 adds defense-in-depth by enforcing the same policy inside the
CodePilot provider itself, so direct provider calls cannot bypass the
\texttt{rasn.tool.agent} gate. Phase 49 adds an optional
\texttt{workspace\_root} sandbox: when configured, read/write/search/list and
replace targets are resolved through rDSN absolute/normalized path helpers and
must remain inside the configured root. Phase 50 adds explicit approval labels:
when \texttt{require\_write\_approval} or
\texttt{require\_shell\_approval} is enabled, CodePilot must attach
\texttt{human\_approved:write} or \texttt{human\_approved:shell} before an
enabled side-effect tool is admitted. Phase 52 adds shell command sandbox
controls: an optional executable allowlist denies shell metacharacters and
requires the first command token to match an allowed executable, while a shell
working-directory wrapper starts commands inside \texttt{workspace\_root} or an
explicit \texttt{shell\_working\_directory}. Tool output is bounded; oversized
output is spilled to an artifact file, and a compact artifact index is written
through \texttt{rasn.state}. A later hardening step adds a configurable shell
executor: \texttt{shell\_executor=container} wraps an approved command through
\texttt{shell\_container\_template}, whose placeholders describe the quoted
command and workspace path. The default executor remains local for compatibility.

This design separates three concerns that are often conflated in prompt-centric
agents: the model decides that a tool is needed, the coordinator routes to an
agent that can perform it, and the policy layer decides whether the particular
side effect is permitted. The current policy engine is intentionally simple, but
the boundary now carries both static configuration decisions and explicit human
approval decisions, leaving room for richer allowlists and audit rules.
For side-effecting tools, the runtime also writes an \texttt{external.effect}
ledger event for policy denials, committed results, replay hits, and replay
misses. This gives operators an effect-level audit stream independent of the
bounded tool-output preview.

Tool execution returns a structured \texttt{tool\_result} with success, output,
and error. Read-like tools are designed for bounded context construction:
\texttt{read} accepts a maximum byte count, \texttt{search} returns matching
lines rather than whole files, and directory listing avoids recursive traversal
unless a future tool explicitly adds it. Write-like tools are disabled by
default because they cross from observation into mutation. When enabled, write
and replace use a temporary file followed by
\texttt{dsn::utils::filesystem::rename\_path} instead of truncating the target
in place. This leverages rDSN filesystem helpers and avoids partially written
target files when a flush or replacement fails. Their outputs still pass through
the same policy and artifact handling paths.

Shell execution has an additional preemption boundary. The configuration key
\texttt{[rasn.policy] shell\_timeout\_ms} bounds command runtime; the Windows
implementation launches the command through native process APIs with explicit
stdout/stderr pipes and calls process termination when the deadline expires.
This does not replace container isolation, but it prevents an approved shell
tool from becoming an unbounded local process. The container executor hook adds
a deployment integration point for Docker, containerd, or a site-specific
sandbox wrapper, but rASN intentionally leaves image policy, mount policy, and
container lifecycle enforcement to future product work.

The artifact spill path is a concrete example of service-bound policy. A large
tool result is not silently truncated into the model prompt. Instead, the tool
returns a bounded preview, writes the full payload under an artifact directory,
and stores metadata in \texttt{rasn.state}. This preserves operator visibility
and gives future workflows a durable handle to the data without overloading the
prompt.

\subsection{Secret Redaction Boundary}

\rasn now treats secret redaction as a shared runtime boundary rather than as a
set of ad hoc string filters in individual tools. The implementation adds a
\texttt{redaction} module configured through \texttt{[rasn.policy]} and the
older \texttt{[rasn.codepilot.tools]} compatibility section. The policy can be
disabled for controlled experiments, but it is enabled by default. It combines
three deterministic mechanisms: exact-value replacement for configured literals,
exact-value replacement for the current values of configured token environment
variables, and syntactic redaction of common credential forms such as bearer
tokens, \texttt{password=...}, \texttt{api\_key: ...}, and JSON key/value
pairs.

The redaction boundary is applied before data reaches persistent or
cross-boundary surfaces. \texttt{event\_log::append} redacts event values before
they are stored in memory or written to JSONL traces. The policy manager redacts
tool output and error text before bounded previews are returned or oversized
payloads are spilled to artifact files. The model service and workflow executor
redact prompts, context entries, provider responses, and provider errors before
invoking providers or recording their outputs. Replay paths also redact values
loaded from older traces before returning them to callers, so replay cannot
re-introduce previously captured credential material.

This mechanism is intentionally deterministic and testable: for a fixed config
snapshot, the same text redacts to the same output, preserving trace comparison
and replay debugging. It is not a substitute for a production secret vault,
process isolation, memory-safe credential handles, or provider SDKs that avoid
temporary token materialization. It does, however, close the immediate prototype
gap in which a read/search/shell result, user prompt, model response, or replay
trace could accidentally persist obvious credential material.

\subsection{State Service}

\texttt{state\_store} validates schema versions and namespaced keys, assigns
monotonic sequence numbers, and supports:

\begin{itemize}
  \item unconditional put,
  \item create-only put,
  \item check-sequence put,
  \item key lookup,
  \item prefix query,
  \item checkpoint,
  \item journal recovery.
\end{itemize}

The workflow subsystem uses check-sequence writes to avoid stale lease updates
and completion races. Policy and observability components write metadata
through injectable state writers that are set to service-graph writers during
graph startup. Checkpoints are written atomically through temporary files and
can be accompanied by append-only journals. The recovery path can optionally
import checkpoint or journal files through rDSN NFS settings before replaying
local state, preserving the same validation logic for local and imported state.
It can also mirror checkpoint and journal files to a local replica directory
through \texttt{[rasn.state.replica]}. The mirror is write-through: if it is
enabled and cannot be updated, the state operation fails explicitly rather than
silently losing the secondary recovery artifact. Recovery can seed missing
primary files from the local replica before attempting NFS import.

The store is small but intentionally strict. A key must be namespaced as
\texttt{scope/id}; raw secret and credential kinds are rejected because the
state service should store secret references rather than secret values. A
create-only write fails if the key already exists. A check-sequence write fails
unless the stored sequence equals the caller's expected sequence. These two
operations are enough to express the current workflow lease and terminal-state
rules.

Checkpoint/recovery is implemented as follows:

\begin{lstlisting}
checkpoint:
  acquire state lock
  copy current records and last sequence
  write temp checkpoint file
  rename temp file into place

recover:
  optionally seed checkpoint/journal from local replica
  optionally import checkpoint/journal through rDSN NFS
  validate checkpoint records
  replay journal records
  rebuild key-value map and last sequence
\end{lstlisting}

The prototype is not a quorum-replicated store, so it does not provide
linearizable recovery across failures. Its value is that it forces all important
runtime state to cross a narrow API and now offers local mirror and NFS seeding
paths for operator recovery. Replacing the backing store with SKV or another
replicated service becomes an implementation substitution rather than an
architectural rewrite.

\subsection{Workflow Compiler and Execution Service}

\texttt{workflow.*} implements declarative workflow parsing and the in-memory
graph executor. A workflow node has an id, action, prompt or tool command,
dependency list, capability, policy labels, timeout budget, retry budget,
latency estimate, cost hint, reliability hint, and optional state key. The
compiler accepts both the compact line-oriented format and structured JSON
workflow specs; both normalize into the same \texttt{workflow\_node} model.
Parsing normalizes metadata, rejects cycles, and builds dependency stages. The optimizer
does not change semantics; it reports a deterministic plan with stage count,
ready-set ordering, maximum parallelism, critical-path latency, cost, and
minimum reliability so operators can inspect tradeoffs before execution.

The current workflow language is intentionally compact. A representative node is:

\begin{lstlisting}
task inspect tool "list ." capability tool.run
  policy read_only budget_ms 5000 latency_ms 50
  cost_hint 1 reliability 98 state unit/inspect
\end{lstlisting}

The same node can be produced by JSON:

\begin{lstlisting}
{
  "nodes": [{
    "id": "inspect",
    "action": "tool",
    "prompt": "list .",
    "capability": "tool.run",
    "policy_labels": ["read_only"],
    "budget_ms": 5000,
    "latency_hint_ms": 50,
    "cost_hint": 1,
    "reliability_hint": 98,
    "state_key": "unit/inspect"
  }]
}
\end{lstlisting}

Model nodes use \texttt{ask}; tool nodes use \texttt{tool}; dependencies are
declared with \texttt{after} in text specs or \texttt{depends\_on} in JSON
specs. JSON support is intentionally strict: malformed types, unsupported
actions, invalid numeric ranges, duplicate ids, invalid state keys, and cycles
fail before any model or tool side effect. The compiler builds a directed
acyclic graph and then computes execution stages. If a node fails, dependent
nodes are marked \texttt{blocked} rather than executed with missing inputs. If a
completed node is available during resume, its output is injected into downstream
context so later model calls can reuse prior work.

The workflow service lifts the graph executor into a durable service. Starting a
run validates and compiles the source, acquires the run lease through a
conditional state write, persists the whole-run \texttt{running} record,
executes nodes through coordinator/model/tool boundaries, persists per-node
progress, and writes a terminal run record. Long-running nodes renew their lease
through an rDSN timer in service mode or a small local renewal thread in direct
mode. Duplicate active starts are rejected; stale leases can be taken over after
the configured TTL.

Resume and cancellation are also state-backed. Resume reloads completed node
records, seeds dependency outputs, marks skipped nodes as \texttt{resumed}, and
executes only incomplete downstream nodes. On startup, the workflow app
schedules delayed recovery after state readiness instead of synchronously
querying state during concurrent \rdsn app startup. On-demand query recovery
loads \texttt{workflow/<run-id>} when the run is not in memory. Phase 46
hardened this path: a genuine \texttt{state key not found} is treated as an
absent run, but any other state failure is propagated. Phase 47 extends the same
semantics to cancellation: \texttt{workflow cancel} recovers a missing
in-memory run from \texttt{rasn.state}, writes the \texttt{cancelled} transition
with a conditional state update, and reports the state error if that durable
transition fails.

The durable execution path is:

\begin{lstlisting}
start(request):
  normalize run id and workflow identity
  recover existing run if needed
  reject duplicate non-resume run
  recover node state for resume run
  acquire execution lease with create-only/CAS state write
  persist run = running
  start lease renewal
  execute graph:
    before each node: check cancellation
    persist node = running
    dispatch ask/tool through shared provider boundary
    persist node = completed/failed/blocked/cancelled
  stop lease renewal
  if cancelled record exists:
    release lease as cancelled
    return cancelled
  persist terminal record with expected running sequence
  release lease as completed/failed
\end{lstlisting}

The conditional terminal write is essential. Without it, a workflow that was
cancelled while a long-running node was still executing could later overwrite the
cancelled state with a completed state. The implementation prevents that by
requiring the terminal completion write to match the sequence of the running
record it started from.

\subsection{Observability and Replay Support}

\texttt{observability.*} provides event queries, failure queries, replay loads,
diagnostic summaries, timelines, and snapshots. Runtime events are structured
records with sequence number, timestamp, trace id, event kind, component, task
id, message, and key-value details. The observability store can answer bounded
queries for interactive use and can persist snapshot indexes through
\texttt{rasn.state}. Snapshot records include trace file, event count, failure
count, and last sequence, giving operators a compact durable pointer into the
larger event stream.

Replay support is intentionally incremental but architecturally important. The
simulator records nondeterministic choices, provider calls record request and
response events, workflow execution records node transitions, and state records
capture durable decisions. Model calls now contribute executable replay data:
prior \texttt{llm.response} events can satisfy provider calls in recorded order.
Tool calls do the same for matching \texttt{tool.ok} and \texttt{tool.error}
events, while missing side-effecting tool results produce an explicit replay
failure rather than repeating the effect. Workflow starts are now executable
replay data as well: the runtime consumes recorded
\texttt{workflow.node.start} events in order and rejects a run whose scheduler
chooses a different next node. File-tool snapshots are also executable replay
data: CodePilot records a \texttt{filesystem.snapshot} event for each
read/list/search target and compares the current fingerprint against the
recorded one before reading local files during replay. A future deterministic
replay engine can combine these artifacts with stronger external-effect
virtualization; the current implementation already gives debugging tools a
stable schema instead of requiring transcript scraping.

The event taxonomy currently includes task lifecycle events, LLM request/response
events, LLM response chunk events, tool request/response events,
nondeterministic choices, retry events, failure events, workflow transitions,
external-effect ledger events, and snapshot events. Query APIs support
bounded event scans, failure-only scans, timelines, diagnostic summaries, replay
loading from trace files, and snapshot persistence. The cursor returned by a
query is the last observed sequence, which lets CLI commands and future UIs build
incremental views.

Observability is intentionally separate from state. State records are compact
facts used for recovery and coordination; observability events are an append-like
history used for diagnosis. A workflow run can be recovered from state without
reading the event log, and a diagnosis can inspect the event log without
mutating state. This separation follows the same principle as distributed-system
logs and checkpoints: recovery state should be minimal and well-structured,
whereas diagnostic traces can be richer and larger.

\subsection{Runtime Metrics}

Structured traces explain a single run; aggregate metrics make a fleet
operable. \texttt{metrics.*} adds a process-global \texttt{metrics\_registry}
backed by \rdsn \texttt{perf\_counter} counters in the \texttt{rasn} section,
reusing the same counter infrastructure \rdsn ships for its own services. The
single choke point \texttt{nucleus\_runtime::record\_event} increments a
cumulative \texttt{COUNTER\_TYPE\_NUMBER} counter for each event kind (Prometheus
\texttt{\_total} convention) and a lazily created counter for each failure class,
so adding a new event kind needs no new call site. Task and model latencies are
paired inside the runtime by task id --- model calls finalize their timing on
both success and failure, so failed-call latency is recorded rather than lost ---
and tool latency is timed at \texttt{rasn\_tool\_agent\_service::run\_tool}, the
point where the CLI, workflow, RPC-server, and direct-call tool paths converge
(an RPC-client transport failure that never reaches the tool agent is the one
exception, and is timed at the service-graph facade instead);
all three feed \texttt{COUNTER\_TYPE\_NUMBER\_PERCENTILES} counters whose
p50/p95/p99/p999 are computed by \rdsn's counter timers.

Two design choices keep this safe to embed everywhere. First, perf-counter
creation asserts a current \rdsn service node, so creation is guarded by a
node-context check and every update is null and lock protected; the metric path
degrades to a no-op (never a crash) in inline, CLI, and node-less contexts, and
\texttt{[rasn.metrics] enabled = false} disables it entirely. Second, the metric
data types and the three renderers (text, Prometheus exposition format, and JSON)
live in a dependency-light translation unit (\texttt{metrics\_format.cpp}) with no
\rdsn includes, so they are unit tested in isolation from the counter backend. A
snapshot is reachable through \codepilot's \texttt{observe metrics
[text|prometheus|json]} and, in service mode, through a \texttt{rasn.metrics}
command registered once with \rdsn's \texttt{command\_manager} and answered over
the local or remote \rdsn CLI.

\subsection{Model Circuit Breaker}

Metrics make a degrading model endpoint observable; the model gateway also reacts
to one. \texttt{circuit\_breaker.*} provides a small, dependency-light engine: a
per-key consecutive-failure breaker with three states (closed, open, half-open)
and a single half-open probe. \texttt{rasn\_llm\_agent\_service} keeps one breaker
per provider and consults it in both \texttt{complete} and
\texttt{complete\_streaming}, at the point where every model path converges on the
provider call. After \texttt{[rasn.model] circuit\_breaker\_failure\_threshold}
consecutive failures the breaker opens and short-circuits requests for
\texttt{circuit\_breaker\_open\_ms}; it then admits one probe that closes the
breaker on success or reopens it on failure. A short-circuited request returns a
normal failed \texttt{llm\_response} carrying the breaker state, so the fast-fail
happens before the provider's own retry loop runs.

The engine mirrors the metrics layer's two design choices. First, it includes no
\rdsn headers and the caller injects the clock as a \texttt{uint64\_t now\_ms}
(\texttt{dsn\_now\_ms}); because that clock is routed through the pluggable
\texttt{env\_provider}, breaker timing is deterministic under replay and the
engine is unit-tested without a live node. Replayed runs and in-process
providers (the deterministic simulator, the workflow service-graph bridge, and
test fakes, identified by \texttt{llm\_provider::in\_process()}) bypass the
breaker, so existing behavior and tests are unchanged; loopback HTTP providers
(Ollama, llama.cpp, LM Studio) are \emph{not} exempt despite a local descriptor,
since they still issue HTTP that can hang. Second, breaker activity is reported
through the same
\texttt{record\_event} choke point as every other metric, adding two counters
(\texttt{rasn\_model\_breaker\_open\_total},
\texttt{rasn\_model\_breaker\_short\_circuit\_total}); live per-provider state is
reachable through \codepilot's \texttt{observe resilience} and a
\texttt{rasn.resilience} command registered with \rdsn's
\texttt{command\_manager}. Tool gateways are not breaker-guarded (they use
admission/rate controls instead); remote-agent RPC dispatch has its own
retryable-failure breaker described below.

\subsection{Model Admission Control}

The circuit breaker stops calling a \emph{broken} provider; admission control keeps
a \emph{healthy} provider from being overwhelmed. A model endpoint that is up but
slow has finite useful concurrency, and pushing past it only deepens queues and
inflates tail latency until the endpoint itself fails (at which point the breaker
takes over). \texttt{admission\_gate.*} therefore fronts each provider with a
dependency-light gate that pairs a concurrency \emph{bulkhead} --- a hard cap of
\texttt{[rasn.model] max\_concurrent\_requests} simultaneous in-flight calls, beyond
which requests are rejected with a normal failed \texttt{llm\_response} --- with
graceful \emph{backpressure}: once in-flight load exceeds
\texttt{soft\_concurrent\_requests}, a short delay (bounded by
\texttt{max\_backpressure\_ms}) is applied before the call so bursts are smoothed
into a steadier arrival rate.

The design reuses \rdsn rather than reinventing it in two ways. First, the
backpressure curve is built from \rdsn's \texttt{exp\_delay} primitive
(\texttt{include/dsn/utility/exp\_delay.h}), the same staged-delay mechanism the host
uses to throttle its task queues, seeded so the delay ramps from zero at the soft
threshold to the configured ceiling; because \texttt{exp\_delay} works entirely in
signed \texttt{int}, the ceiling is clamped to an int-safe bound before it scales the
curve, so an out-of-range \texttt{max\_backpressure\_ms} degrades to a
large-but-valid delay rather than overflowing negative. Second, the gate reuses the
breaker's \texttt{in\_process()} predicate to exempt no-network providers, and is
wired into \texttt{rasn\_llm\_agent\_service} on the same convergence point as the
breaker. Ordering is deliberate: the admission \emph{reservation} is evaluated
\emph{before} the breaker's authoritative probe-consuming check (so a rejected
request never perturbs breaker state) and \emph{after}
the replay check (so replayed runs stay deterministic), while the graceful
backpressure \emph{delay} is applied only after the breaker also admits --- so a
request the breaker short-circuits neither sleeps nor holds its slot through a
sleep. An RAII \texttt{admission\_slot} releases capacity on every exit, including
an early breaker short-circuit. With the generous
defaults, sequential single-agent CLI use (in-flight $\approx 1$) never delays or
rejects; the gate engages only under the concurrent fan-out of service mode.
Rejections and backpressure delays are exported as
\texttt{rasn\_model\_admission\_rejected\_total} and
\texttt{rasn\_model\_admission\_delayed\_total} through the same
\texttt{record\_event} choke point, and live per-provider in-flight depth is shown by
the \texttt{rasn.resilience} command. The cap is per-provider; a single
process-wide concurrency budget is noted as future work.

\subsection{Model Rate Limiter}

The breaker reacts to a \emph{broken} provider and admission control protects a
\emph{healthy} one from concurrency overload; the rate limiter governs the third
classic dimension, \emph{throughput}. Hosted model APIs publish
requests-per-minute and tokens-per-minute quotas, and exceeding them does not
degrade gracefully: the endpoint returns HTTP~429s that burn the retry budget
and, on metered services, cost money. A concurrency bulkhead does not prevent
this, because a single caller issuing rapid sequential requests stays at
in-flight $\approx 1$ yet can still exceed a per-minute quota.
\texttt{rate\_limiter.*} therefore fronts each provider with a token-bucket
governor: the bucket starts full (absorbing an initial burst up to its capacity),
refills at the configured sustained rate, and on each call either admits
immediately, \emph{paces} the request with a short delay that reserves the next
token, or --- when the projected wait would exceed
\texttt{[rasn.model] rate\_limit\_max\_wait\_ms} --- rejects it with a normal failed
\texttt{llm\_response} instead of blocking a worker.

The engine is dependency-light like the breaker and admission gate: its header
pulls in no \rdsn or thrift types, and the current time is \emph{injected} as a
millisecond value rather than read from a clock. The caller passes
\texttt{::dsn\_now\_ms()}, so token refill flows through \rdsn's pluggable
environment provider --- making it deterministic under replay and unit-testable
without a live node, exactly as the breaker's cooldown clock is. Refill is hardened
against a non-monotonic injected clock: a reading at or before the last refill
timestamp neither removes tokens nor moves the refill baseline backward, since
moving it back would let a later-but-still-stale reading measure a positive interval
and refill prematurely; the baseline holds at its high-water mark until the clock
advances past it. It is wired into
\texttt{rasn\_llm\_agent\_service} beside the other two gates. Ordering is again
deliberate. A non-mutating open-breaker precheck runs \emph{first}: a request to a
provider whose breaker is already open (and still cooling down) fast-fails ahead of
any admission or rate rejection, so a broken dependency wins over an overload or
quota signal without wasting an admission slot or rate token; the precheck only
reads breaker state and never consumes the one-shot half-open probe. Otherwise
admission reserves a slot and the rate limiter acquires a token first, and the
breaker's authoritative check --- which consumes a one-shot half-open probe that
must be paired with a result report --- is consulted \emph{last}, so a request that
admission or the rate limiter rejects never strands a half-open probe. If the
breaker's authoritative check does short-circuit a request that already acquired a
rate token --- the half-open-probe-busy case, or a race where the breaker opened
between the precheck and the authoritative check --- that token is \emph{refunded}
before the fast-fail returns (and the admission slot is released by its RAII guard),
so a breaker short-circuit never silently drains the provider's quota or delays its
recovery probe; the refund is a no-op when the limiter is disabled and is clamped to
burst capacity. The pacing
delay, like
admission backpressure, is applied only after the breaker admits, and the two
delays are coalesced into a single wait (the larger of the two) so they never
stack into additive over-delay. The same \texttt{in\_process()} predicate exempts
no-network providers. The limiter is disabled by default
(\texttt{rate\_limit\_requests\_per\_min}~$=0$ means unlimited), so existing
single-agent CLI and deterministic test behavior is unchanged; an operator opts in
when pointing \rasn at a metered endpoint. Rejections and pacing delays are
exported as \texttt{rasn\_model\_rate\_limited\_total} and
\texttt{rasn\_model\_rate\_delayed\_total} through the same \texttt{record\_event}
choke point, and live per-provider rate, burst, and available tokens are shown by
the \texttt{rasn.resilience} command. Rate limiting governs request count;
token/cost budgets and a process-wide budget are noted as future work.

\subsection{Tool Gateway Admission and Rate Controls}

The tool gateway is another dependency boundary: local filesystem tools, shell
commands, and future remote tool adapters can be fanned out by workflows or model
responses independently of model-provider health. Earlier versions protected tools
with policy, replay, workspace scoping, and timeouts, but they did not provide an
explicit overload governor. \rasn now reuses \texttt{admission\_gate.*} and
\texttt{rate\_limiter.*} for tool execution, keyed by tool name.

The ordering preserves existing safety semantics. \texttt{rasn\_tool\_agent\_service}
first validates the generic request, checks for a recorded replay result, fails
closed on missing side-effect replay, and evaluates \texttt{rasn.policy}. Only an
actual, policy-approved provider call reserves a tool admission slot or rate token.
Thus a replayed shell/write result is deterministic and a policy-denied tool call
does not consume capacity. When both gates admit, their requested delays are
coalesced into a single wait (the larger of admission backpressure and rate pacing)
before invoking the provider. The \texttt{admission\_slot} RAII guard releases
capacity on success, tool failure, output-bound truncation, validation failure after
execution, or any early return.

The implementation also removes accidental serialization from the provider pointer
lock. The service now takes a \texttt{std::shared\_ptr} to the current
\texttt{agent\_tool\_provider} while holding \rdsn's \texttt{zlock}, releases the
lock, and then runs validation, policy, governors, and provider execution with that
stable reference. Concurrency is therefore controlled by product-level
\texttt{[rasn.tool]} settings rather than by an all-tool mutex. The new events
\texttt{tool.admission.rejected}, \texttt{tool.admission.delayed},
\texttt{tool.rate.limited}, and \texttt{tool.rate.delayed} feed four counters and
the expanded \texttt{rasn.resilience} command, which includes model and tool
gateway state. Tool admission defaults are generous for sequential \codepilot usage,
and the tool rate limiter remains unlimited unless an operator sets
\texttt{rate\_limit\_requests\_per\_min}.

\subsection{Remote-Agent Dispatch Resilience}

The coordinator's service-mode RPC dispatch path is also an outbound dependency.
Registry lookup may select a built-in model/tool service or a future custom agent,
but the coordinator still needs to protect itself from an unhealthy or slow remote
agent. \rasn therefore adds a per-agent guard around
\texttt{coordinator\_router::invoke\_remote}: the key is the selected descriptor's
\texttt{agent\_id}, and the guarded operation is the outbound
\texttt{RPC\_RASN\_AGENT\_INVOKE}.

The guard reuses the same engines as the model gateway. Endpoint preflight
(\texttt{coordinator\_router::validate\_remote\_endpoint}) resolves the descriptor
address once and rejects malformed or unresolvable endpoints before any breaker
probe, admission slot, or rate token is consumed, emitting a counted
\texttt{remote\_agent.endpoint.invalid} event; the resolved address is then reused
for dispatch to avoid a second resolution. A non-mutating breaker
precheck runs first, then a per-agent admission slot and rate token are acquired,
then the authoritative breaker check admits or short-circuits the call. If the
breaker short-circuits after a rate token has been acquired, the token is refunded;
admission and rate delays are coalesced into one sleep before the RPC. This
preserves the model-gateway invariant that admission/rate rejections never strand a
half-open breaker probe. The guard is installed only on the RPC-client branch of
\texttt{rasn\_coordinator\_service::invoke}, so inline standalone execution is not
changed.

The breaker report deliberately distinguishes dependency failures from
application outcomes. \texttt{coordinator\_router::invoke\_remote} marks transport
failures as retryable; remote agents may also return retryable transient errors.
Only those retryable outcomes are reported as breaker failures. Non-retryable
responses --- policy denials, validation failures, deterministic tool errors, or
other application-level decisions --- count as successful dependency responses so a
valid but rejected request does not poison the remote-agent circuit breaker.

The configuration lives under \texttt{[rasn.remote\_agent]}:
\texttt{circuit\_breaker\_*}, \texttt{admission\_enabled},
\texttt{max\_concurrent\_requests}, \texttt{soft\_concurrent\_requests},
\texttt{max\_backpressure\_ms}, and \texttt{rate\_limit\_*}. Defaults are generous
(\texttt{64} hard concurrency, \texttt{32} soft concurrency, \texttt{200ms}
backpressure, rate unlimited) to preserve normal service-smoke behavior. Runtime
events feed seven counters:
\texttt{rasn\_remote\_agent\_breaker\_open\_total},
\texttt{rasn\_remote\_agent\_breaker\_short\_circuit\_total},
\texttt{rasn\_remote\_agent\_admission\_rejected\_total},
\texttt{rasn\_remote\_agent\_admission\_delayed\_total},
\texttt{rasn\_remote\_agent\_rate\_limited\_total},
\texttt{rasn\_remote\_agent\_rate\_delayed\_total}, and
\texttt{rasn\_remote\_agent\_endpoint\_invalid\_total}; live state appears in
\texttt{observe resilience} and the \texttt{rasn.resilience} command alongside the
model and tool sections.

\subsection{Testing and Validation}

The current unit suite covers message validation, endpoint metadata, registry
leases, coordinator retries, policy classification, artifact state writes,
observability snapshot state writes, service-graph lifecycle, workflow parsing,
workflow cancellation, restart-safe cancel recovery, cancel persistence failure
reporting, budget propagation, resume, state checkpoint/recovery, conditional
writes, workflow execution leases, lease renewal, cancellation tombstones,
state-recovery error propagation, metric label sanitization, metric snapshot
rendering, cumulative counter deltas from runtime events, per-failure-class
counters, remote-agent resilience counters, and the \texttt{rasn.metrics} and
\texttt{rasn.resilience} commands. The CodePilot tool tests cover
provider-side default-deny write/shell policy and atomic file replacement. Direct
\codepilot self-tests
exercise model health, generic model invocation, tool description, state
put/get/checkpoint, workflow compile/run/resume/cancel, and observability
snapshot. Service-mode smoke starts the full \rdsn app graph and runs the same
self-test through RPC.

\section{\codepilot: A Coding CLI on \rasn}

\subsection{User Model}

\codepilot is a command-line application inspired by existing coding agents.
It supports one-shot prompts, an interactive REPL, explicit tool commands,
workflow commands, registry inspection, agent control, observability commands,
skills, provider configuration, trace/replay, and a self-test. The CLI is not
the runtime; it is an application adapter over \rasn.

Table~\ref{tab:codepilot-commands} groups the current command surface. The
important design point is that most commands exercise a generic \rasn service
surface rather than a CodePilot-only implementation path.

\begin{table}[t]
\caption{\codepilot command groups.}
\label{tab:codepilot-commands}
\centering
\begin{tabular}{ll}
\toprule
Group & Purpose \\
\midrule
\texttt{ask}/\texttt{stream}/prompt & Model completion and streamed chunks through coordinator \\
\texttt{tool} & Explicit local tool execution \\
\texttt{workflow} & Validate, compile, start, resume, query, cancel, nodes \\
\texttt{state} & Put, get, query, checkpoint/recover smoke \\
\texttt{registry} & Inspect descriptors and capability routing \\
\texttt{agentctl} & Describe, heartbeat, query, cancel agent requests \\
\texttt{observe} & Events, failures, timelines, snapshots, replay, runtime metrics \\
\texttt{schema} & Text, JSON, IDL, C++ contracts/clients, TypeScript, and Python export \\
\texttt{eval} & Built-in and external-CLI benchmark harness \\
\texttt{selftest} & End-to-end direct/service regression workload \\
\bottomrule
\end{tabular}
\end{table}

\subsection{Command Routing}

For model requests, \codepilot builds an \texttt{agent\_completion\_request},
which is converted into a generic \texttt{agent\_request} with capability
\texttt{model.complete}. For tools, it builds an \texttt{agent\_tool\_request}
or a generic \texttt{agent\_request} with capability \texttt{tool.run}. These
requests are sent through \texttt{rasn\_service\_graph}, which calls the
coordinator inline or through RPC depending on mode.

\codepilot has three execution shapes. In direct mode, a command owns a
command-scoped service-graph reference and calls the in-process graph. In
service mode, \codepilot is an rDSN app in a larger app graph and talks to
service endpoints through clientlets. In self-test mode, it executes a fixed
end-to-end workload that covers provider health, model invocation, tool
description, state put/get/checkpoint, workflow compile/run/resume/cancel,
observability snapshot, and service readiness. This makes \codepilot both a
usable example application and a regression driver for the generic runtime.

Prompt construction is deliberately conservative. Context files are read as
explicit context entries, not silently concatenated into hidden global state.
Timeouts and retry budgets are carried in the runtime request. Skill prompts are
application-level templates that produce ordinary model requests. This keeps the
runtime independent of CodePilot's prompt style and makes it possible for other
agent applications to reuse the same services with different prompts.

\subsection{Readiness and Error Reporting}

Before dispatching service-mode commands, \codepilot probes every critical
surface: state, registry, coordinator, model-agent identity, model health,
tool-agent identity, workflow validation, and observability. A readiness failure
names the missing boundary rather than returning a generic CLI failure. Command
errors preserve structured runtime errors where possible, including
state-recovery failures, lease conflicts, model health failures, policy denials,
and cooperative cancellation results.

\subsection{Local Tools}

\codepilot provides a tool provider with:

\begin{itemize}
  \item \texttt{list <path>},
  \item \texttt{read <path> [max-bytes]},
  \item \texttt{search <path> <text>},
  \item \texttt{write <path> <content>},
  \item \texttt{replace <path> <old> <new>},
  \item \texttt{shell <command>}.
\end{itemize}

Read/search/list are read-only, normalize platform separators, and require
existing targets. Write and shell tools are disabled by default and require
explicit configuration. Enabled write/replace/shell calls can additionally
require human approval: direct CLI users can pass \texttt{tool --yes}, while
agent-requested tool calls trigger an interactive prompt before CodePilot adds a
side-effect-specific approval label. Tool output is bounded; oversized output is
written to an artifact directory and indexed in state. Side-effecting tool
intents are also recorded as \texttt{external.effect} events, and approved shell
commands can run either locally or through a configured container command
template.

The local tool provider is intentionally small because it is a policy boundary,
not a full shell automation framework. \texttt{list} is for fast orientation,
\texttt{read} is for bounded context import, \texttt{search} is for targeted
evidence gathering, \texttt{write} and \texttt{replace} are controlled mutation
primitives, and \texttt{shell} is a last-resort escape hatch. In a production
CodePilot, these would be complemented by diff preview, richer audit review, and
hardened process or container orchestration for command execution.

\subsection{Skills}

Skills are \codepilot-specific reusable prompt templates, not generic \rasn
runtime primitives. Current built-ins include rDSN plugin work, code review,
build debugging, feature planning, and documentation. They demonstrate how an
application can layer domain-specific prompt assets on top of the generic
runtime without polluting core service abstractions.

\subsection{Provider Modes}

The simulator provider enables offline testing and deterministic replay
experiments. HTTP providers connect to Copilot-compatible, OpenAI-compatible,
Ollama, llama.cpp, and LM Studio endpoints. This design supports both cloud and
local model serving, but the current provider layer is intentionally minimal
and should be hardened before production use.

\subsection{Configuration and Examples}

\codepilot is configured through normal \rdsn configuration files plus
provider-specific sections. Example files under \texttt{examples/} cover direct
workflow execution, service-mode smoke, provider selection, state recovery, and
tool-policy behavior. This is intentionally close to the way rDSN applications
are normally deployed: the CLI can be used as a single local binary during
development, but the same code can participate in an rDSN app graph when the
runtime services need to be inspected independently.

The examples also serve as executable documentation for the platform. A workflow
example demonstrates the declarative graph; a service smoke configuration starts
the full rDSN app graph; provider examples show local and HTTP-compatible model
configuration; and self-test exercises the minimum viable product surface. This
is important for research reproducibility because future experiments can pin a
configuration file, provider mode, and workflow input rather than relying on an
interactive transcript.

\section{Experimental Design}
\label{sec:experimental-design}

\rasn is not yet finished enough for a definitive empirical evaluation.
Nevertheless, the prototype now includes an executable evaluation harness so
future experiments can reuse a stable procedure instead of a prose-only plan.
\texttt{codepilot eval} runs built-in or file-backed task suites through
\codepilot and reports success, latency, and output size. \texttt{codepilot eval
external <template>} invokes an external CLI command template containing
\texttt{\{prompt\}}, enabling side-by-side comparisons with OpenCode-, Codex-,
or Copilot-style command-line agents under the same task prompts.

\subsection{Research Questions}

\begin{description}
  \item[RQ1: Reliability.] Does \rasn reduce unrecoverable or ambiguous agent
  failures compared with transcript-centric coding CLIs?
  \item[RQ2: Diagnosability.] Do \rasn traces, state records, and observability
  queries help developers localize failures faster?
  \item[RQ3: Reproducibility.] Can captured nondeterminism and workflow state
  support replay or partial replay of failed tasks?
  \item[RQ4: Overhead.] What latency, CPU, memory, and storage overhead does
  the service graph introduce in inline and \rdsn service modes?
  \item[RQ5: Generality.] Can the same runtime support coding, review,
  planning, and non-coding agent workflows without changing core services?
\end{description}

\subsection{Environment Setup}

The primary development environment is Windows with Visual Studio 2022 and
\rdsn's Windows build scripts. Cross-platform validation runs in WSL with a
full TeX Live installation for this report and a Linux \rdsn build. The
evaluation should include:

\begin{itemize}
  \item Windows 11, Visual Studio 2022, local \rdsn checkout, Boost 1.84.
  \item Ubuntu WSL or native Ubuntu with the same \rdsn revision.
  \item Local simulator provider for deterministic and offline tests.
  \item Local model providers through Ollama, llama.cpp, or LM Studio.
  \item Cloud provider adapters for OpenAI-compatible or Copilot-compatible
  endpoints when credentials and policy permit.
  \item Identical task repositories and benchmark scripts across all tools.
\end{itemize}

\subsection{Workloads}

The workload suite should contain one prompt per non-empty, non-comment line
when used with \texttt{codepilot eval <suite-file>}. It should cover:

\begin{enumerate}
  \item \textbf{Microbenchmarks}: model call, tool call, state put/get, workflow
  compile, workflow start, registry query, observability query, cancellation,
  and recovery.
  \item \textbf{Coding tasks}: bug fixes, refactorings, test additions,
  documentation updates, and build-failure diagnosis across small and medium
  repositories.
  \item \textbf{Failure-injection tasks}: state service unavailability, model
  timeout, tool denial, duplicate workflow start, stale lease takeover,
  cancellation storm, oversized tool output, and invalid workflow source.
  \item \textbf{Replay tasks}: simulator traces and selected model traces where
  nondeterministic choices can be captured and replayed.
\end{enumerate}

\subsection{Metrics}

The evaluation should report:

\begin{itemize}
  \item task success rate and patch correctness,
  \item number of failed or blocked workflow nodes,
  \item mean and tail latency,
  \item token/provider cost where available,
  \item number of tool calls and denied side-effect attempts,
  \item trace completeness and number of observable failure records,
  \item recovery success after restart,
  \item duplicate-execution prevention rate under lease contention,
  \item cancellation acknowledgement and terminal-cancellation rate,
  \item storage overhead of traces, state records, journals, and artifacts,
  \item developer diagnosis time in controlled debugging studies.
\end{itemize}

\subsection{Comparison Targets}

The comparison should include:

\begin{itemize}
  \item \textbf{\codepilot inline mode}: baseline \rasn without RPC boundaries.
  \item \textbf{\codepilot service mode}: full \rdsn service graph.
  \item \textbf{OpenCode-style CLI}: representative open-source coding CLI.
  \item \textbf{Codex-style CLI}: a modern model-centric coding CLI.
  \item \textbf{Claude Code or similar}: commercial coding-agent CLI when
  accessible.
  \item \textbf{GitHub Copilot CLI-style workflow}: comparison against a
  developer-facing Copilot command-line experience where permitted.
\end{itemize}

Because commercial systems change rapidly and may not expose all internals,
comparisons should focus on externally observable behavior: success, latency,
failure reporting, reproducibility support, and operator controls.

\subsection{Evaluation Procedure}

Each task should be run multiple times with fixed repository snapshots,
recorded prompts, provider settings, and random seeds where possible. For
model-dependent tasks, the evaluation should separate provider quality from
runtime behavior by using both the simulator and real models. Experiments can
begin with:

\begin{lstlisting}[language=bash]
codepilot eval tasks.txt
codepilot eval external "opencode run --prompt {prompt}" tasks.txt
codepilot eval external "codex exec {prompt}" tasks.txt
\end{lstlisting}

The harness reports per-task success, latency, output bytes, and aggregate
failure counts. External commands are templates rather than hard-coded
integrations because CLI syntax and access policies change quickly.
Failure-injection
experiments should verify that errors are surfaced at the intended component
boundary rather than converted into misleading success-shaped responses.

\section{Discussion}

\subsection{Threats to Validity}

\textbf{Internal validity.} Runtime improvements may appear to improve task
success when they only improve failure visibility. Evaluation must separate
correctness, diagnosability, and overhead. Provider variability can dominate
results, so experiments need fixed model versions where possible.

\textbf{External validity.} The first application is a coding CLI. Results may
not directly generalize to browser agents, data-analysis agents, or long-lived
autonomous services. However, the core runtime concerns--state, RPC boundaries,
observability, leases, cancellation, and nondeterminism--are common across many
agent domains.

\textbf{Construct validity.} Metrics such as "diagnosability" are hard to
measure. The evaluation should combine quantitative trace completeness with
controlled human studies or structured debugging tasks.

\textbf{Implementation validity.} The current implementation is a prototype.
Some components are intentionally lightweight, including the curl provider,
single-node state service, and same-process service-mode smoke. These choices
are useful for research iteration but not sufficient for production claims.

\subsection{Practicality}

\rasn's practical value is strongest where operators need robust execution
rather than only fluent interaction. Coding agents in enterprise repositories
need audit trails, policy gates, provider configurability, checkpointing, and
recovery. \rasn provides early versions of these capabilities while preserving a
simple local CLI path.

The main practical cost is complexity. A service graph, state service,
workflow engine, and observability layer are heavier than a direct script that
calls an LLM. The design therefore supports inline mode for local prototyping
and service mode for integration and robustness testing.

\subsection{Generality}

\rasn is not inherently tied to coding. The generic protocol uses capabilities,
tasks, context entries, and tool providers. Future applications could implement
research assistants, data-analysis agents, incident-response agents, or
multi-agent review systems. The main requirement is that application-specific
prompts and tools remain outside the core runtime.

\subsection{Extensibility}

The most important extension points are:

\begin{itemize}
  \item new LLM providers,
  \item new tool providers,
  \item external agent descriptors in the registry,
  \item stronger state backends,
  \item richer policy engines,
  \item alternative workflow optimizers,
  \item additional observability exporters,
  \item production deployment topologies using \rdsn URI resolution.
\end{itemize}

\subsection{Remaining Product Gaps}

Table~\ref{tab:remaining-gaps} summarizes the current product limitations after
the latest hardening round. These are no longer prototype blockers, but they are
the main differences between the current \rasn implementation and a
production-complete platform.

\begin{table*}[t]
\caption{Remaining product limitations after the current hardening round.}
\label{tab:remaining-gaps}
\centering
\begin{tabular}{p{0.16\textwidth}p{0.38\textwidth}p{0.36\textwidth}}
\toprule
Area & Current capability & Remaining limitation \\
\midrule
State availability & Checkpoints, append-only journals, conditional writes,
workflow leases, optional local replica mirroring, and optional \rdsn NFS import
& No quorum-replicated or externally managed HA state backend yet \\
Tool isolation & Default-deny side effects, workspace scoping, approvals,
command allowlists, timeout/job containment, and a configurable container
command wrapper & No hardened container orchestrator with image, mount,
network, and lifecycle policy \\
Replay fidelity & Model-response replay, matching tool-result replay,
workflow-schedule replay, filesystem snapshots, and an
\texttt{external.effect} ledger for side-effect intents & No full virtualization
of arbitrary external services, clocks, network state, or process environments \\
Deployment validation & Inline mode, typed service-mode RPC, URI/host endpoint
configuration, registry heartbeats, and active lease cleanup & Multi-process and
cluster deployment tests remain limited \\
Credential storage & \texttt{token\_ref} handles for environment variables,
files, and commands plus deterministic redaction & No vault-backed or OS-backed
credential provider integration \\
SDK packaging & Generated C++/TypeScript/Python contracts and RPC-client source
& Packaged SDKs and concrete TypeScript/Python transports are not shipped \\
Evaluation evidence & Unit tests, self-tests, service smokes, schema smokes,
report build, and a small eval harness & Large benchmarks and user studies for
debugging effectiveness remain future work \\
Observability metrics & Cumulative event counters and task/model/tool latency
percentiles over \rdsn \texttt{perf\_counter}, exported as text, Prometheus, and
JSON via \texttt{observe metrics} and the \texttt{rasn.metrics} command & No
bundled scrape gateway, retention store, or prebuilt dashboards; latency
percentiles depend on \rdsn's periodic counter timers \\
Overload/dependency isolation & Model providers and remote-agent RPC dispatch
have circuit-breaker, admission, and rate controls; tools have per-tool admission
and rate controls; all are surfaced through \texttt{observe resilience} and
\texttt{rasn.resilience} & Limits are per-dependency rather than process-wide and
request-count based rather than token/cost based \\
\bottomrule
\end{tabular}
\end{table*}

\section{Related Work}

\textbf{Distributed systems and service runtimes.} \rdsn is the direct
foundation for this work. More broadly, \rasn draws on distributed-systems
ideas including explicit service identity, typed RPC, leases, checkpointing,
failure detectors, and observable execution. Classical work on replicated state machines, Paxos, and distributed
snapshots~\cite{lamport1998parttime,chandy1985distributed} motivates the
emphasis on state, ownership, and traceability, although the current prototype
is not a replicated system.

\textbf{Workflow and orchestration systems.} Workflow engines and dataflow
systems have long exposed task graphs, retries, checkpoints, and recovery.
\rasn applies similar ideas to agent workflows, where nodes can be model calls
or tool calls and nondeterminism is a first-class concern.

\textbf{LLM agent frameworks.} Recent systems such as ReAct, Toolformer,
Reflexion, AutoGen, LangChain, and related frameworks demonstrate how LLMs can
reason, call tools, coordinate multiple agents, and reflect on
failures~\cite{yao2023react,schick2023toolformer,shinn2023reflexion,wu2023autogen,langchain}.
\rasn is complementary: it focuses less on prompting strategy and more on the
runtime substrate for robustness, service boundaries, state, and observability.

\textbf{Coding agents.} \codepilot is inspired by command-line coding agents
such as OpenCode~\cite{opencode}, Codex-style CLIs, Claude Code, and GitHub
Copilot CLI-style interfaces. The distinction is that \codepilot is designed
as a sample application over a generic service nucleus rather than as a
standalone model-tool loop.

\textbf{Observability and replay.} Debugging nondeterministic systems has a
long history in software engineering. \rasn's current replay support is modest,
but its event model is intended to evolve toward stronger diagnosis and replay
tools for agent behavior.

\section{Conclusion}

\rasn explores a systems-oriented path toward robust AI agent platforms. By
building on \rdsn, the prototype treats agent components as distributed
services with explicit lifecycle, typed RPC, state, leases, cancellation,
readiness, and observability. \codepilot demonstrates how a coding CLI can be
implemented as an application adapter over this nucleus while keeping generic
runtime concerns reusable. The system remains a prototype, and a rigorous
empirical evaluation is future work. Nevertheless, the current implementation
shows that many agent robustness concerns can be expressed as concrete runtime
abstractions and incrementally hardened through tests, service-mode smoke
checks, and documented invariants.

\bibliographystyle{ACM-Reference-Format}
\bibliography{main}

\end{document}
